\documentclass[10pt,conference]{IEEEtran}
\usepackage{amsmath,amssymb,stmaryrd}
\usepackage[T1]{fontenc}
\usepackage{graphicx}
\usepackage{booktabs}
\usepackage{array}
\usepackage{tabularx}
\usepackage{adjustbox}
\usepackage[table]{xcolor}
\usepackage{pifont}
\newcommand{\yes}{\ding{51}}\newcommand{\no}{\ding{55}}
\usepackage{tikz}
% Harvey balls for the comparison table: does, partly does, does not.
\DeclareRobustCommand{\hb}[1]{\tikz[baseline=-0.55ex]{\draw[line width=0.4pt] (0,0) circle (0.75ex);
  \ifcase#1\or\fill (0,0.75ex) arc (90:270:0.75ex) -- cycle;\or\fill (0,0) circle (0.75ex);\fi}}
\DeclareRobustCommand{\full}{\hb{2}}\DeclareRobustCommand{\half}{\hb{1}}\DeclareRobustCommand{\none}{\hb{0}}\newcommand{\na}{--}
\usepackage{url}
\usepackage{algorithm}
\usepackage[noend]{algpseudocode}
\providecommand{\ket}[1]{\lvert #1\rangle}
\newcommand{\Bw}{B_{\mathrm{QDG}}}
\newcommand{\Fw}{F_{\mathrm{QDG}}}
% The tool's name, in one place. tools/build_anon.sh defines \anonbuild for the double-anonymous submission.
\ifdefined\anonbuild\newcommand{\tool}{QLine}\else\newcommand{\tool}{srcSlice-Q}\fi
\newcommand{\TODO}[1]{\textcolor{red}{[TODO: #1]}}
\newtheorem{definition}{Definition}
\newtheorem{theorem}{Theorem}
\newtheorem{proposition}{Proposition}
% Generated by experiments/icpc/make_numbers.py from committed results. Do not edit by hand.

\newcommand{\numNPairs}{784}  % results/icpc/rq1/records.csv
\newcommand{\numNCircuits}{158}  % results/icpc/rq1/records.csv
\newcommand{\numRqOneDuUnsound}{216}  % results/icpc/rq1/records.csv
\newcommand{\numRqOneDuUnsoundPct}{27.6\%}  % results/icpc/rq1/records.csv
\newcommand{\numRqOneDuRecall}{0.718}  % results/icpc/rq1/records.csv
\newcommand{\numRqOneDuPrecision}{0.970}  % results/icpc/rq1/records.csv
\newcommand{\numRqOneDuSize}{51.2\%}  % results/icpc/rq1/records.csv
\newcommand{\numRqOneLcUnsound}{0}  % results/icpc/rq1/records.csv
\newcommand{\numRqOneLcUnsoundPct}{0.0\%}  % results/icpc/rq1/records.csv
\newcommand{\numRqOneLcRecall}{1.000}  % results/icpc/rq1/records.csv
\newcommand{\numRqOneLcPrecision}{0.923}  % results/icpc/rq1/records.csv
\newcommand{\numRqOneLcSize}{70.4\%}  % results/icpc/rq1/records.csv
\newcommand{\numRqOneYzUnsound}{0}  % results/icpc/rq1/records.csv
\newcommand{\numRqOneYzUnsoundPct}{0.0\%}  % results/icpc/rq1/records.csv
\newcommand{\numRqOneYzRecall}{1.000}  % results/icpc/rq1/records.csv
\newcommand{\numRqOneYzPrecision}{0.835}  % results/icpc/rq1/records.csv
\newcommand{\numRqOneYzSize}{78.6\%}  % results/icpc/rq1/records.csv
\newcommand{\numRqOneQUnsound}{0}  % results/icpc/rq1/records.csv
\newcommand{\numRqOneQUnsoundPct}{0.0\%}  % results/icpc/rq1/records.csv
\newcommand{\numRqOneQRecall}{0.807}  % results/icpc/rq1/records.csv
\newcommand{\numRqOneQPrecision}{0.972}  % results/icpc/rq1/records.csv
\newcommand{\numRqOneQSize}{54.0\%}  % results/icpc/rq1/records.csv
\newcommand{\numRqOneFourEdgeUnfixedUnsound}{5}  % results/icpc/rq1/records.csv
\newcommand{\numRqOneFourEdgeUnfixedUnsoundPct}{0.6\%}  % results/icpc/rq1/records.csv
\newcommand{\numRqOneFourEdgeUnfixedRecall}{0.994}  % results/icpc/rq1/records.csv
\newcommand{\numRqOneFourEdgeUnfixedPrecision}{0.864}  % results/icpc/rq1/records.csv
\newcommand{\numRqOneFourEdgeUnfixedSize}{75.0\%}  % results/icpc/rq1/records.csv
\newcommand{\numRqOneFourEdgeFixedUnsound}{0}  % results/icpc/rq1/records.csv
\newcommand{\numRqOneFourEdgeFixedUnsoundPct}{0.0\%}  % results/icpc/rq1/records.csv
\newcommand{\numRqOneFourEdgeFixedRecall}{1.000}  % results/icpc/rq1/records.csv
\newcommand{\numRqOneFourEdgeFixedPrecision}{0.864}  % results/icpc/rq1/records.csv
\newcommand{\numRqOneFourEdgeFixedSize}{75.2\%}  % results/icpc/rq1/records.csv
\newcommand{\numRqOneDuMissKickback}{1{,}737}  % results/icpc/rq1/misses.csv
\newcommand{\numRqOneDuMissSameWire}{835}  % results/icpc/rq1/misses.csv
\newcommand{\numRqOneDuMissWrittenOperand}{654}  % results/icpc/rq1/misses.csv
\newcommand{\numRqOneDuMissFeedback}{134}  % results/icpc/rq1/misses.csv
\newcommand{\numRqOneQRemoved}{2{,}847}  % results/icpc/rq1/pairs.csv
\newcommand{\numRqOneQRemovedNonInfluencing}{669}  % results/icpc/rq1/pairs.csv
\newcommand{\numRqOneQRemovedContributing}{1{,}618}  % results/icpc/rq1/pairs.csv
\newcommand{\numRqOneQPrunedPairs}{477}  % results/icpc/rq1/pairs.csv
\newcommand{\numLcqApplicable}{732}  % results/icpc/rq1_lcq_remap/records.csv
\newcommand{\numLcqNotApplicable}{52}  % results/icpc/rq1_lcq_remap/records.csv
\newcommand{\numLcqNaMeasurements}{42}  % results/icpc/rq1_lcq_remap/records.csv
\newcommand{\numLcqUnsound}{0}  % results/icpc/rq1_lcq_remap/records.csv
\newcommand{\numLcqSmallerThanLc}{243}  % results/icpc/rq1_lcq_remap/records.csv
\newcommand{\numLcqRecall}{0.93}  % results/icpc/rq1_lcq_remap/records.csv
\newcommand{\numQVsLcqReduction}{12.3\%}  % results/icpc/rq1_lcq_remap/records.csv
\newcommand{\numQVsLcqReductionPooled}{10.8\%}  % results/icpc/rq1_lcq_remap/records.csv
\newcommand{\numQLargerThanLcq}{72}  % results/icpc/rq1_lcq_remap/records.csv
\newcommand{\numHOnePHolm}{\ensuremath{p < 10^{-10}}}  % results/icpc/confirmatory/summary.json
\newcommand{\numHTwoDiff}{16.3}  % results/icpc/confirmatory/summary.json
\newcommand{\numHTwoCILow}{15.1}  % results/icpc/confirmatory/summary.json
\newcommand{\numHTwoCIHigh}{17.7}  % results/icpc/confirmatory/summary.json
\newcommand{\numHTwoRelative}{23.2\%}  % results/icpc/confirmatory/summary.json
\newcommand{\numHTwoSmaller}{477}  % results/icpc/confirmatory/summary.json
\newcommand{\numHTwoPHolm}{\ensuremath{p < 10^{-10}}}  % results/icpc/confirmatory/summary.json
\newcommand{\numHTwoR}{1.00}  % results/icpc/confirmatory/summary.json
\newcommand{\numHThreeRate}{55.4\%}  % results/icpc/confirmatory/summary.json
\newcommand{\numHThreeCILow}{51.9\%}  % results/icpc/confirmatory/summary.json
\newcommand{\numHThreeCIHigh}{58.8\%}  % results/icpc/confirmatory/summary.json
\newcommand{\numHThreePrunedRate}{26.6\%}  % results/icpc/confirmatory/summary.json
\newcommand{\numHThreePrunedRobust}{127}  % results/icpc/confirmatory/summary.json
\newcommand{\numHThreePrunedN}{477}  % results/icpc/confirmatory/summary.json
\newcommand{\numHFourContainQ}{82.2\%}  % results/icpc/confirmatory/summary.json
\newcommand{\numHFourContainLc}{95.7\%}  % results/icpc/confirmatory/summary.json
\newcommand{\numHFourDiff}{\ensuremath{-}13.5}  % results/icpc/confirmatory/summary.json
\newcommand{\numHFourCILow}{\ensuremath{-}14.4}  % results/icpc/confirmatory/summary.json
\newcommand{\numHFourCIHigh}{\ensuremath{-}12.6}  % results/icpc/confirmatory/summary.json
\newcommand{\numHFiveSizeDiff}{2.9}  % results/icpc/confirmatory/summary.json
\newcommand{\numHFiveExamDiff}{\ensuremath{-}0.061}  % results/icpc/confirmatory/summary.json
\newcommand{\numRobQPrunedBasis}{65.8\%}  % results/icpc/rq2_robustness/pairs.csv
\newcommand{\numRobLcBasis}{100.0\%}  % results/icpc/rq2_robustness/pairs.csv
\newcommand{\numRobQPrunedHaar}{26.6\%}  % results/icpc/rq2_robustness/pairs.csv
\newcommand{\numRobLcHaar}{100.0\%}  % results/icpc/rq2_robustness/pairs.csv
\newcommand{\numRobQPrunedParam}{78.1\%}  % results/icpc/rq2_robustness/pairs.csv
\newcommand{\numRobLcParam}{100.0\%}  % results/icpc/rq2_robustness/pairs.csv
\newcommand{\numRobLcqBasisRobust}{732}  % results/icpc/rq2_robustness_lcq/pairs.csv
\newcommand{\numRobLcqBasisN}{732}  % results/icpc/rq2_robustness_lcq/pairs.csv
\newcommand{\numRobLcqHaarRobust}{732}  % results/icpc/rq2_robustness_lcq/pairs.csv
\newcommand{\numRobLcqHaarN}{732}  % results/icpc/rq2_robustness_lcq/pairs.csv
\newcommand{\numRobLcqParamRobust}{171}  % results/icpc/rq2_robustness_lcq/pairs.csv
\newcommand{\numRobLcqParamN}{175}  % results/icpc/rq2_robustness_lcq/pairs.csv
\newcommand{\numAblUnfixedFull}{0.5\%}  % results/rq4_ablation/TAB_ablation_rerun.csv
\newcommand{\numAblUnfixedNoED}{54.0\%}  % results/rq4_ablation/TAB_ablation_rerun.csv
\newcommand{\numAblUnfixedNoMD}{13.4\%}  % results/rq4_ablation/TAB_ablation_rerun.csv
\newcommand{\numAblUnfixedNoUED}{13.3\%}  % results/rq4_ablation/TAB_ablation_rerun.csv
\newcommand{\numAblUnfixedNoCD}{4.8\%}  % results/rq4_ablation/TAB_ablation_rerun.csv
\newcommand{\numAblFixedFull}{0.0\%}  % results/icpc/ablation_fixed/TAB_ablation_rerun.csv
\newcommand{\numAblFixedNoED}{63.9\%}  % results/icpc/ablation_fixed/TAB_ablation_rerun.csv
\newcommand{\numAblFixedNoMD}{13.0\%}  % results/icpc/ablation_fixed/TAB_ablation_rerun.csv
\newcommand{\numAblFixedNoUED}{26.1\%}  % results/icpc/ablation_fixed/TAB_ablation_rerun.csv
\newcommand{\numAblFixedNoCD}{6.7\%}  % results/icpc/ablation_fixed/TAB_ablation_rerun.csv
\newcommand{\numAblPairs}{962}  % results/icpc/ablation_fixed/TAB_ablation_rerun.csv
\newcommand{\numFamCircuits}{119}  % results/icpc/families/rows.csv
\newcommand{\numFamLcUnsound}{0}  % results/icpc/families/rows.csv
\newcommand{\numFamYzUnsound}{0}  % results/icpc/families/rows.csv
\newcommand{\numFamQUnsound}{0}  % results/icpc/families/rows.csv
\newcommand{\numFamDuUnsound}{232}  % results/icpc/families/rows.csv
\newcommand{\numFamQSoundNotMinimal}{212}  % results/icpc/families/rows.csv
\newcommand{\numFamQRows}{819}  % results/icpc/families/rows.csv
\newcommand{\numCatEntries}{24}  % results/icpc/catalogue/catalogue.json
\newcommand{\numRqThreeUnits}{5{,}688}  % results/icpc/rq3/records.csv
\newcommand{\numRqThreeContainDu}{77.3\%}  % results/icpc/rq3/records.csv
\newcommand{\numRqThreeContainLc}{95.7\%}  % results/icpc/rq3/records.csv
\newcommand{\numRqThreeContainYz}{97.4\%}  % results/icpc/rq3/records.csv
\newcommand{\numRqThreeContainQ}{82.2\%}  % results/icpc/rq3/records.csv
\newcommand{\numRqThreeContainRandom}{62.4\%}  % results/icpc/rq3/records.csv
\newcommand{\numCQFour}{769}  % results/icpc/rq3/units.csv
\newcommand{\numCQFourRecovered}{691}  % results/icpc/rq3/units.csv
\newcommand{\numBugsFourQKilled}{7}  % results/icpc/rq3_bugs4q/bugs.csv
\newcommand{\numBugsFourQInjured}{1}  % results/icpc/rq3_bugs4q/bugs.csv
\newcommand{\numBugsFourQAlive}{1}  % results/icpc/rq3_bugs4q/bugs.csv
\newcommand{\numBugsFourQExcluded}{2}  % results/icpc/rq3_bugs4q/bugs.csv
\newcommand{\numBugsFourQTriaged}{107}  % bugs4q/triage.csv
\newcommand{\numBugsFourQIncluded}{11}  % bugs4q/triage.csv
\newcommand{\numMapInstructions}{137}  % results/icpc/source_map/bugs.csv
\newcommand{\numMapMapped}{136}  % results/icpc/source_map/bugs.csv
\newcommand{\numMapCrossCompared}{127}  % results/icpc/source_map/crosscheck.csv
\newcommand{\numMapCrossAgree}{127}  % results/icpc/source_map/crosscheck.csv
\newcommand{\numMapCrossPrograms}{5}  % results/icpc/source_map/crosscheck.csv
\newcommand{\numLinesFullContained}{15}  % results/icpc/source_map/units.csv
\newcommand{\numLinesFullMean}{10.7}  % results/icpc/source_map/units.csv
\newcommand{\numLinesLcContained}{15}  % results/icpc/source_map/units.csv
\newcommand{\numLinesLcMean}{10.3}  % results/icpc/source_map/units.csv
\newcommand{\numLinesQContained}{15}  % results/icpc/source_map/units.csv
\newcommand{\numLinesQMean}{9.5}  % results/icpc/source_map/units.csv
\newcommand{\numLinesDuContained}{14}  % results/icpc/source_map/units.csv
\newcommand{\numLinesDuMean}{4.9}  % results/icpc/source_map/units.csv
\newcommand{\numLinesUnits}{15}  % results/icpc/source_map/units.csv
\newcommand{\numMbflUnits}{293}  % results/icpc/mbfl/units.csv
\newcommand{\numMbflExamAll}{0.348}  % results/icpc/mbfl/units.csv
\newcommand{\numMbflExamQ}{0.465}  % results/icpc/mbfl/units.csv
\newcommand{\numMbflExamLc}{0.419}  % results/icpc/mbfl/units.csv
\newcommand{\numTTwoSonnetFull}{84.4\%}  % results/llm/t2_scores.csv
\newcommand{\numTTwoSonnetLc}{77.2\%}  % results/llm/t2_scores.csv
\newcommand{\numTTwoSonnetQ}{69.9\%}  % results/llm/t2_scores.csv
\newcommand{\numTTwoSonnetDu}{62.3\%}  % results/llm/t2_scores.csv
\newcommand{\numTTwoSonnetQVsFullDiff}{\ensuremath{-}14.5}  % results/llm/t2_scores.csv
\newcommand{\numTTwoSonnetQVsFullCI}{[\ensuremath{-}19.9, \ensuremath{-}9.4]}  % results/llm/t2_scores.csv
\newcommand{\numTTwoSonnetQVsFullP}{\ensuremath{p = 1.6\times 10^{-7}}}  % results/llm/t2_scores.csv
\newcommand{\numTTwoSonnetQVsLcDiff}{\ensuremath{-}7.2}  % results/llm/t2_scores.csv
\newcommand{\numTTwoSonnetQVsLcCI}{[\ensuremath{-}12.3, \ensuremath{-}2.5]}  % results/llm/t2_scores.csv
\newcommand{\numTTwoSonnetQVsLcP}{\ensuremath{p = 0.006}}  % results/llm/t2_scores.csv
\newcommand{\numTTwoHaikuFull}{71.4\%}  % results/llm/t2_scores.csv
\newcommand{\numTTwoHaikuLc}{65.9\%}  % results/llm/t2_scores.csv
\newcommand{\numTTwoHaikuQ}{58.7\%}  % results/llm/t2_scores.csv
\newcommand{\numTTwoHaikuDu}{51.8\%}  % results/llm/t2_scores.csv
\newcommand{\numTTwoHaikuQVsFullDiff}{\ensuremath{-}12.7}  % results/llm/t2_scores.csv
\newcommand{\numTTwoHaikuQVsFullCI}{[\ensuremath{-}18.8, \ensuremath{-}6.5]}  % results/llm/t2_scores.csv
\newcommand{\numTTwoHaikuQVsFullP}{\ensuremath{p = 8.2\times 10^{-5}}}  % results/llm/t2_scores.csv
\newcommand{\numTTwoHaikuQVsLcDiff}{\ensuremath{-}7.2}  % results/llm/t2_scores.csv
\newcommand{\numTTwoHaikuQVsLcCI}{[\ensuremath{-}12.7, \ensuremath{-}2.2]}  % results/llm/t2_scores.csv
\newcommand{\numTTwoHaikuQVsLcP}{\ensuremath{p = 0.010}}  % results/llm/t2_scores.csv
\newcommand{\numTTwoQwenMidFull}{42.8\%}  % results/llm/t2_scores.csv
\newcommand{\numTTwoQwenMidLc}{43.1\%}  % results/llm/t2_scores.csv
\newcommand{\numTTwoQwenMidQ}{36.2\%}  % results/llm/t2_scores.csv
\newcommand{\numTTwoQwenMidDu}{33.3\%}  % results/llm/t2_scores.csv
\newcommand{\numTTwoQwenMidQVsFullDiff}{\ensuremath{-}6.5}  % results/llm/t2_scores.csv
\newcommand{\numTTwoQwenMidQVsFullCI}{[\ensuremath{-}11.6, \ensuremath{-}1.4]}  % results/llm/t2_scores.csv
\newcommand{\numTTwoQwenMidQVsFullP}{\ensuremath{p = 0.020}}  % results/llm/t2_scores.csv
\newcommand{\numTTwoQwenMidQVsLcDiff}{\ensuremath{-}6.9}  % results/llm/t2_scores.csv
\newcommand{\numTTwoQwenMidQVsLcCI}{[\ensuremath{-}11.2, \ensuremath{-}2.5]}  % results/llm/t2_scores.csv
\newcommand{\numTTwoQwenMidQVsLcP}{\ensuremath{p = 0.003}}  % results/llm/t2_scores.csv
\newcommand{\numTTwoQwenSmallFull}{13.0\%}  % results/llm/t2_scores.csv
\newcommand{\numTTwoQwenSmallLc}{12.7\%}  % results/llm/t2_scores.csv
\newcommand{\numTTwoQwenSmallQ}{13.0\%}  % results/llm/t2_scores.csv
\newcommand{\numTTwoQwenSmallDu}{13.4\%}  % results/llm/t2_scores.csv
\newcommand{\numTTwoQwenSmallQVsFullDiff}{0.0}  % results/llm/t2_scores.csv
\newcommand{\numTTwoQwenSmallQVsFullCI}{[\ensuremath{-}4.3, 4.3]}  % results/llm/t2_scores.csv
\newcommand{\numTTwoQwenSmallQVsFullP}{\ensuremath{p = 1.00}}  % results/llm/t2_scores.csv
\newcommand{\numTTwoQwenSmallQVsLcDiff}{0.4}  % results/llm/t2_scores.csv
\newcommand{\numTTwoQwenSmallQVsLcCI}{[\ensuremath{-}2.9, 3.6]}  % results/llm/t2_scores.csv
\newcommand{\numTTwoQwenSmallQVsLcP}{\ensuremath{p = 1.00}}  % results/llm/t2_scores.csv
\newcommand{\numTTwoItems}{276}  % results/llm/t2_scores.csv
\newcommand{\numCostLcMaxMs}{16}  % results/icpc/cost/circuits.csv
\newcommand{\numCostQDepthTen}{1.4}  % results/icpc/cost/circuits.csv
\newcommand{\numCostBuildDepthTen}{0.03}  % results/icpc/cost/circuits.csv
\newcommand{\numCostQDepthTwenty}{10.6}  % results/icpc/cost/circuits.csv
\newcommand{\numCostBuildDepthTwenty}{0.27}  % results/icpc/cost/circuits.csv
\newcommand{\numCostQDepthForty}{49.9}  % results/icpc/cost/circuits.csv
\newcommand{\numCostBuildDepthForty}{1.10}  % results/icpc/cost/circuits.csv
\newcommand{\numCostDepthEightyOverLimit}{4}  % results/icpc/cost/circuits.csv
\newcommand{\numCostBuildDepthEighty}{5.2}  % results/icpc/cost/circuits.csv
\newcommand{\numCostDepthOneSixtyOverLimit}{5}  % results/icpc/cost/circuits.csv
\newcommand{\numCostBuildDepthOneSixty}{21.7}  % results/icpc/cost/circuits.csv
\newcommand{\numCostWallLimit}{120}  % icpc/cost_run.py
\newcommand{\numCostWidthQMax}{8.7}  % results/icpc/cost/circuits.csv
\newcommand{\numCorpusLCircuits}{55}  % results/icpc/corpus_l/circuits.csv
\newcommand{\numCorpusLCompleted}{39}  % results/icpc/corpus_l/circuits.csv
\newcommand{\numNoiseLcCrossExactMax}{\ensuremath{4.2\times 10^{-15}}}  % results/icpc/noise/pairs.csv
\newcommand{\numNoiseLcCrossExactNonzero}{0}  % results/icpc/noise/pairs.csv
\newcommand{\numNoiseLcCrossExactN}{1{,}746}  % results/icpc/noise/pairs.csv
\newcommand{\numNoiseLcCrossShotMax}{0.033}  % results/icpc/noise/pairs.csv
\newcommand{\numNoiseQCrossExactMax}{\ensuremath{2.4\times 10^{-1}}}  % results/icpc/noise/pairs.csv
\newcommand{\numNoiseQCrossExactNonzero}{702}  % results/icpc/noise/pairs.csv
\newcommand{\numNoiseQCrossExactN}{1{,}746}  % results/icpc/noise/pairs.csv
\newcommand{\numNoiseQCrossShotMax}{0.033}  % results/icpc/noise/pairs.csv
\newcommand{\numNoiseShots}{1{,}024}  % icpc/noise_run.py
\newcommand{\numHwCircuits}{18}  % results/icpc/hardware/hardware_tvd_corrected.csv
\newcommand{\numHwWithin}{16}  % results/icpc/hardware/hardware_tvd_corrected.csv
\newcommand{\numHwNoiseFreeZero}{18}  % results/icpc/hardware/hardware_tvd_corrected.csv


\begin{document}
\title{What Must a Quantum Slice Keep? Slicing Quantum Programs for Comprehension}
\author{\IEEEauthorblockN{Anonymous Author(s)}}
\maketitle

\begin{abstract}
Which statements must a slice of a quantum program keep to preserve the state of selected qubits? The dependence rules
proposed for quantum programs have not been evaluated against this question, which is fundamental to using slicing for
program comprehension. We define a correct slice as one that preserves the selected qubits' joint state for every input
state, and a dependence as an operation after which a change can alter that state. Our approach builds the quantum
dependence graph (QDG) from a write-all rule, under which a gate writes every qubit it acts on, its controls included.
We prove that QDG slices are correct and retain every dependence, and that classical data, classical control and
quantum data dependence suffice, so further kinds proposed for quantum programs, such as entanglement dependence, add
no dependence the QDG misses. We compare the QDG with alternative rules on \numNPairs{} circuit--qubit pairs and
\numStudyFaultUnits{} fault units. Treating controls as read-only, as the classical program dependence graph (PDG) does,
gives incorrect slices for \numCpdgUnsoundPct{} of the pairs, and \numFdCpdgZMatchPct{} of these hide the error because
the qubit's measurement outcomes match those of the whole program. The other proposed quantum PDG that we compare with
keeps \numYzLcExcessRatio{} times as many operations that cannot affect the selected qubits as the QDG, and
pruning slices to make them smaller removes genuine dependences. QDG slices contain \numRqThreeContainLc{} of the seeded
faults and take milliseconds to compute, but they do not improve the ranking of mutation-based fault localization.
These findings provide a semantic foundation for quantum slicing and clarify its benefits and limits for
comprehension and debugging.
\end{abstract}

\begin{IEEEkeywords}
program slicing, program comprehension, quantum programs, dependence analysis, light cone, empirical study
\end{IEEEkeywords}

% Sections are written one at a time and agreed with the author before the next starts.
\section{Introduction}\label{sec:intro}

A developer who sees a qubit end in the wrong state asks which lines of the program determine it. For classical
programs a backward slice answers this~\cite{weiser1984program}, and programmers use slices to read and debug
code~\cite{weiser1982programmers, gallagher1991using}. Quantum programs need such help more. Their intermediate states
cannot be inspected without disturbing them, and a gate can change a qubit it does not target, as in phase kickback,
where a controlled gate changes the phase of its control without entangling the two
qubits~\cite{nielsen2010quantum}. Quantum developers debug
by printing, drawing circuits, simulating and reading code, and report that this does not
scale~\cite{zappin2026challenges}. Bug detectors such as LintQ and QChecker find known patterns by static analysis,
which reads a program without running it, % static-definition
but many faults are wrong outputs that no pattern describes~\cite{paltenghi2024lintq, zhao2023qchecker}.

Slicing has been proposed for quantum programs~\cite{yu2025qpdg} and named as an aid to understanding
them~\cite{murillo2025roadmap}, but no proposed dependence rule has been checked against what a quantum program computes. For comprehension this matters
most, because a slice that hides code able to change the qubit sends its reader away from the fault, and such slices
look right. The classical slice of the textbook program of Fig.~\ref{fig:bv} reproduces the faulty program's state
and leaves out the faulty line. Over the \numNPairs{} circuit--qubit pairs of our study, reading a controlled gate as
an assignment that leaves its control unchanged gives a wrong slice for \numCpdgUnsoundPct{} of the pairs, and
\numFdCpdgZMatchPct{} of those reproduce the whole program's outcome histogram, so a developer who checks a slice by
its outcomes would accept them.

This paper asks what a quantum slice must keep, and which dependence rule produces it. A slice of a set of qubits must
preserve their joint state for every input, and an operation is a dependence of them if a change right after it can
change that state. One rule meets both, the \emph{write-all rule}, under which a gate writes every qubit it touches,
its controls included. On qubits its slice is the backward light cone, whose soundness compilers already rely
on~\cite{murali2019triq, qiskit2025}. What has been missing is a proof that it keeps every dependence, so that no
other proposed kind of dependence can add one, and a measurement of what each alternative loses. We give the proof and
measure the losses on \numNCircuits{} benchmark circuits with an exact ground truth, \numFamCircuits{} generated
circuits, \numRqThreeUnits{} seeded-fault units and \numBugsFourQUnits{} real-fault units from
Bugs4Q~\cite{zhao2023bugs4q}. The study was pre-registered, but its confirmatory tests concern pruning, and the
comparisons of the other rules, added after the first results, are exploratory.

The paper makes four contributions. The first is a definition of what a slice of a set of qubits must preserve and of
dependence, which separates sound slices from slices that keep every dependence. The second is a proof that the slice
of the quantum dependence graph (QDG), defined by the write-all rule and extended along classical bits, is sound for
every input, with the induction mechanized in Isabelle/HOL, and keeps every dependence, with evidence that none of Qiskit's
standard two-qubit gates can safely write fewer qubits. The third shows that classical data, classical control and quantum data
dependence suffice, since apparent quantum control dependence is quantum data dependence and entanglement is neither
necessary nor sufficient for dependence. The fourth measures the alternatives. The classical PDG fails almost
entirely through its reading of the control. The QPDG of Yu and Zhao~\cite{yu2025qpdg} keeps \numYzLcExcessRatio{}
times as much code that cannot matter as the QDG, whose slices are on average \numLcExcess{} percentage points larger
than the smallest slice that keeps every dependence, a gap that grows with width. Pruning drops dependences. Slices
contain most faults but improve neither mutation-based fault ranking nor a language model's localization of faults.

The paper builds no slicer for a particular language. Its definitions are stated on the circuit a program builds.
Section~\ref{sec:lines} shows how a slice reaches source lines, and computing slices at scale on the source developers
write is a separate question. Section~\ref{sec:related} reviews related work, Section~\ref{sec:example} works through
Fig.~\ref{fig:bv}, Section~\ref{sec:core} defines the semantics and the rules, Sections~\ref{sec:design}
and~\ref{sec:results} present the study, and the last three sections give the implications, the threats to validity
and the conclusion. The data, scripts, protocol and proofs are in an anonymous replication package.\TODO{anonymous package link}


\section{A Classical Slice That Looks Right}\label{sec:example}

\begin{figure}[t]
  \centering
  \includegraphics[width=\columnwidth]{figures/fig_bv.pdf}
  \caption{The Bernstein--Vazirani program of the Qiskit textbook~\cite{qiskittextbook}, with comments and blank lines
  omitted and the textbook's line numbers kept. Line~25 carries a seeded fault, \texttt{x(n)} where the textbook
  applies \texttt{z(n)}. The two columns mark the lines in the backward slice of $q_0$ under the classical PDG and
  under the QDG. The classical slice misses the orange lines, which can change $q_0$. Slices are computed statically
  by \tool{}, and outcomes and dependences by the study's code.}
  \label{fig:bv}
\end{figure}

This section shows, on a textbook program, how a slice computed the classical way can reproduce the output a
developer sees and still hide the fault, and why a check of measured outcomes cannot tell. It starts with the
quantum concepts the rest of the paper uses.

\paragraph*{Background}
A qubit's state is a superposition $\alpha\ket{0} + \beta\ket{1}$, and measuring it yields $0$ or $1$ with
probabilities $|\alpha|^2$ and $|\beta|^2$. The two states $\ket{+}$ and $\ket{-}$ give both outcomes equally often
and differ only in the sign of $\beta$, their phase, which no measurement in the computational basis can see. Gates are
reversible operations on qubits, such as the Hadamard gate H, which turns $\ket{0}$ into $\ket{+}$ and $\ket{+}$ back
into $\ket{0}$. A controlled gate such as the controlled-NOT (CNOT) applies an operation to its target depending on
its control, but it can also change the control. When the target is in $\ket{-}$, the CNOT leaves the target as it is
and flips the phase of its control instead, an effect called phase kickback. Two qubits are entangled when their
joint state cannot be written as a state of each qubit alone. A measurement writes a classical bit, which later gates
may read as a condition. Qiskit programs build circuits in Python, with loops, functions and conditions that compute
which gates act on which qubits.

\paragraph*{The program}
The Bernstein--Vazirani algorithm~\cite{bernstein1997qct} finds a hidden bit string $s$ with a single query to an
oracle. Fig.~\ref{fig:bv} shows the program of the Qiskit textbook~\cite{qiskittextbook}, with $s = 011$. It rests on
phase kickback. Lines~24 and~25 prepare the auxiliary qubit $q_3$ in $\ket{-}$, line~29 puts every data qubit in
$\ket{+}$, and for every bit of $s$ that is $1$, the CNOT of line~40 kicks a phase onto its data qubit, turning
$\ket{+}$ into $\ket{-}$. The Hadamard gates of line~47 turn these phases into the bits of $s$, so the program
returns~$011$. We seeded one fault, a one-token change on line~25 that applies X instead of Z. The auxiliary is
then left in $\ket{+}$, on which a CNOT changes nothing, and the program returns \numBvBuggyOut{} instead of
\numBvFixedOut{}.

\paragraph*{The slices}
The developer sees $q_0$ read $0$ where $s$ says $1$, and slices backward from $q_0$ to see which lines determine it.
A classical slicer reads \texttt{cx(q, n)}, which flips its target when its control is $1$, as the assignment
$q_n \leftarrow q_n \oplus q$, which reads $q$ and writes only the auxiliary. Its slice of $q_0$ keeps the Hadamard
gates of lines~\numBvCpdgLines{} and the lines that compute their qubits, lines \numBvCpdgClassical{}, as the left
column of Fig.~\ref{fig:bv} shows. The slice looks right. Its \numBvCpdgOps{} operations give $q_0$ the same state as
the faulty program, with fidelity \numBvCpdgFidelity{}, a similarity of quantum states that is $1$ exactly when they
are equal, so it gives the developer no reason to look at lines~\numBvMissedLines{}, where the fault is.

These lines can nonetheless change $q_0$. Restoring Z on line~25 turns $q_0$ from $\ket{0}$ into $\ket{1}$, two states
with fidelity \numBvFixFidelity{}, because the CNOT acts on its control as well as its target. A slice that lets
every gate write all its qubits, the QDG of Section~\ref{sec:core}, keeps lines~\numBvQdgLines{}, together with the
lines that compute them, among them the oracle's loop and condition on lines~35 to~39. Its slice contains the fault. The QPDG of Yu and Zhao keeps the same lines here, since its two extra operations, the Hadamard gate and the
CNOT of $q_1$, are later iterations of lines~29 and~40 that cannot change $q_0$, and Fig.~\ref{fig:ansatz} shows where
it costs more.

\paragraph*{Why the outcomes cannot tell}
Here the fault shows in the measured outcome of $q_0$. That is not always so. In a real bug from
Bugs4Q~\cite{zhao2023bugs4q}, posted on Stack Overflow, a two-qubit program meant to prepare
$(\ket{01}+\ket{10})/\sqrt{2}$ applies \texttt{h(0)}, \texttt{h(1)} and \texttt{cx(0,1)}, with \texttt{h(1)} where
\texttt{x(1)} was needed, and returns all four bit strings equally often. Each qubit alone reads $1$ with probability
\numMotPOne{}, in the buggy program and in the intended one alike, so the outcomes of either qubit cannot reveal the
failure, and only their joint outcomes do. The fault does change the state of $q_0$, which is a uniformly random bit in
the intended program and the pure $\ket{+}$ in the buggy one, two states with fidelity \numMotFidelity{}. Because the
CNOT writes its control, the QDG slice of $q_0$ contains the faulty \texttt{h(1)}, while the classical slice keeps only
\texttt{h(0)}.

Both lessons hold beyond these programs, as Section~\ref{sec:results} shows. Reading a gate as an assignment misses
influence that flows into a control qubit, and judging a slice by the outcomes of one qubit misses errors in phase and
in correlations. A slice must therefore be judged by the state of the qubits it is about, which is how
Section~\ref{sec:core} defines it.


\section{Related Work}\label{sec:related}

This section places our work among six lines of research. For each line it states what others have done, how
our work relates to it and what we reuse. It closes with the gaps that remain open and with the data each
line was evaluated on, and Table~\ref{tab:related} summarizes both.

\subsection{Program slicing and srcSlice}
Weiser introduced program slicing~\cite{weiser1984program} and showed that programmers use slices when they
debug~\cite{weiser1982programmers}. Dependence graphs reduced slicing to graph
reachability~\cite{ottenstein1984pdg, ferrante1987pdg, horwitz1990interprocedural}, Podgurski and Clarke gave
dependence a semantic definition~\cite{podgurski1990formal}, and dynamic slicing restricts a slice to one
execution~\cite{korel1988dynamic, agrawal1990dynamic}. Surveys catalog the many kinds of slices that
followed~\cite{tip1994a, silva2012vocabulary, binkley2004a}. Probabilistic slicing preserves the distribution of the
criterion variables~\cite{hur2014slicing, amtoft2020theory, olmedo2025slicing}, and Gallagher names quantum
programs as a direction for semantically defined slicing~\cite{gallagher2025program}.

srcSlice computes slices from srcML, the XML form of source code that keeps every character and position of
the original~\cite{collard2011lightweight, collard2013srcml}. Instead of building a dependence graph, it records
for every variable a \emph{slice profile}, the lines where the variable is defined and used, the variables it flows
into and the calls it is passed to, and it computes dependences from these profiles as it slices. It slices the
Linux kernel in about 20 minutes with accuracy comparable to CodeSurfer~\cite{alomari2014srcslice}. General
frameworks such as WALA provide dependence-graph slicers for Python~\cite{wala}, but on a Qiskit program every
gate is a method call on one circuit object, so neither they nor srcSlice can tell which qubits an operation
depends on.

\emph{Relation and reuse.} We keep the classical notions, a criterion, a backward and a forward slice, and
dependence as the possibility of change, and define them for the reduced state of a set of qubits, a criterion
stronger than a distribution in one basis, which misses phase. The \tool{} approach reuses srcML's
representation and srcSlice's slice profiles for the classical part of a program and adds a quantum layer,
which Section~\ref{sec:approach} describes. We fix two gaps in srcSlice's handling of Python, subscript uses and
module globals, as a patch to the open-source tool.

\subsection{srcML for quantum programs}
Behler et al.\ specify a srcML markup for OpenQASM~3, so that quantum syntax shares one representation with the
classical languages srcML already supports~\cite{behler2024supporting}. On it they build an execution-order
analysis that unrolls loops, expands gate and function calls and labels operations inside \texttt{if}
statements with their condition, and their QStatic tool uses the resulting per-qubit order to detect
iteration and recursion patterns and to suggest refactorings~\cite{behler2025static}. Their analyzer does not
yet handle \texttt{while} loops or qubits passed to functions as classical values~\cite{behler2024supporting},
and its order table marks the control of a CNOT apart from its target. Behler et al.\ name, as future work, the
adaptation of srcML's slicer to find every line that affects a given qubit~\cite{behler2025static}.

\emph{Relation and reuse.} Our work takes that step. We reuse the idea that a static execution order of
qubit operations, kept on srcML, is the basis of quantum analyses, and we add what a slicer needs on top of it,
a dependence rule with a guarantee. An order table serves pattern detection well, but read directly as a
dependence graph it treats the control as read-only, which Section~\ref{sec:core} shows to be unsound. Their
analyzer reads OpenQASM~3 and ours reads Qiskit and OpenQASM~2, so the two cover different languages.

\subsection{Dependence graphs for quantum programs}
Yu and Zhao's quantum program dependence graph (QPDG)~\cite{yu2025qpdg} extends the classical
PDG~\cite{ferrante1987pdg} with entanglement, no-cloning, temporal and measurement dependences and proposes
slicing on it. It is defined in prose with examples and has not been implemented or evaluated. Its own example
shows why validation matters. In the graph Yu and Zhao draw for their first listing, the slice of a measured
qubit omits the two gates that fix its outcome, so the slice returns 1 with probability
\numYzFigSliceProb{} where the program returns 1 with certainty. Kaul et al.\ build a code property graph for
Qiskit and OpenQASM programs and query it for bugs and smells, among them superfluous operations, gates that
do not affect a measured qubit~\cite{kaul2023uniform}. Their graph treats the control of a CNOT as unchanged,
so the query can report gates whose effect reaches a measured qubit back through a control. Kaul et al.\
demonstrate the graph on example programs and note the lack of large datasets.
Roadmaps of quantum software engineering describe dependence graphs as an aid to understanding
programs~\cite{murillo2025roadmap}, and a roadmap of testing and debugging finds dependence and data-flow
analysis of circuits almost unexplored, naming only Kaul et al.~\cite{ramalho2025road}.

\emph{Relation and reuse.} We take the QPDG's question, which operations a qubit depends on, and answer it
against a semantic definition. We reimplement the QPDG as a documented baseline, logging every reading of its
prose before any run, and Section~\ref{sec:core} shows that its entanglement edges add no dependence the
write-all rule misses.

\subsection{Entanglement analysis}
Static analyses approximate which qubits may be entangled, with type systems~\cite{perdrix2007types,
prost2009reasoning, yuan2022twist}, abstract interpretation~\cite{perdrix2008entanglement, yu2021abstract}, the
stabilizer formalism~\cite{honda2015stabilizer} and data-flow rules~\cite{xia2025entanglement}. ScaffCC assumes
that qubits that interact may be entangled~\cite{javadiabhari2015scaffcc}, the rule our QPDG baseline uses.

\emph{Relation.} These analyses answer a different question, whether two qubits may be correlated, and they are
not slicers. Section~\ref{sec:core} shows that entanglement neither implies nor is implied by dependence, so
their results cannot decide what a slice keeps.

\subsection{Light cones and compilers}
Compilers build gate-dependence graphs to schedule and map circuits~\cite{murali2019triq, amy2020staq}, and
QAOA analyses bound the gates that can reach an observable by its light cone~\cite{farhi2014qaoa}. Qiskit's
\texttt{LightCone} pass keeps the gates that can reach an observable and drops those that commute with
it~\cite{qiskit2025}, and Heim et al.\ note that quantum dependences should account for
commutation~\cite{heim2020quantum}. Quantum separation logics reason locally about the qubits a program
touches~\cite{zhou2021qsl, le2022qsl}. Compilers also move a single-qubit gate within the idle time of its qubit
without changing the computation~\cite{smith2022timestitch}, and call the time windows of a circuit its
slices, a sense different from ours.

\emph{Relation and reuse.} The backward light cone is the relation of our write-all rule, and we do not claim
the relation as new. What these works do not provide is a definition of what a slice must preserve and a proof
that the light cone is exactly what a slice must keep, which Section~\ref{sec:core} gives. We reuse Qiskit's
\texttt{LightCone} pass as the commutation-pruning baseline.

\subsection{Debugging tools and bug detectors}
Metwalli and Van Meter adapt slicing to circuits by splitting a circuit at barriers the developer places and
removing the qubits a part does not use, and their tool maps gates to source lines through tracebacks at run
time~\cite{metwalli2022debugging, metwalli2024testing}. Their phase tests detect a wrong phase but need
hand-designed test vectors to locate the faulty qubit. Qiskit's slicing utilities partition a circuit by
barriers, depth or gate type~\cite{qiskitaddonutils}. A roadmap of testing and debugging lists slicing as a way
around no-cloning and notes that it grows hard to manage as circuits grow~\cite{ramalho2025road}, and testing
faces the related limit that phase faults escape computational-basis measurement~\cite{oldfield2024faster}.
Bug detectors find known patterns statically. LintQ runs CodeQL queries over Qiskit programs and orders two
operations when they share a qubit~\cite{paltenghi2024lintq}. QChecker matches patterns in the abstract syntax
tree~\cite{zhao2023qchecker}, and Q-PAC also detects known patterns~\cite{nayak2023qpac}. Bug
patterns~\cite{zhao2021bugpatterns}, Bugs4Q~\cite{zhao2023bugs4q}, mutation tools~\cite{mendiluze2021muskit,
fortunato2022qmutpy, garcia2022testing} and mutation-based fault localization~\cite{ishimoto2025mbfl} supply
faults and rankings to this field.

\emph{Relation and reuse.} Slices complement these tools. Manual slicing has no criterion, so the developer still
narrows the fault by hand, and a slice by the write-all rule generalizes the removal of unused qubits, since it
drops every qubit and operation that cannot reach the criterion. The \tool{} approach maps operations to source
lines statically. LintQ's ordering relation is the relation of the write-all rule, which confirms the rule in
practice, while LintQ uses it for patterns and makes no claim about slices. Detectors report known patterns,
and of the 42 Bugs4Q bugs on which QChecker was evaluated, 22 are wrong outputs that it cannot
detect~\cite{zhao2023qchecker}. A slice shows where such a fault can be. We reuse Bugs4Q, LintQ's corpus and its
public warnings, published mutation operators and mutation-based fault localization.

\subsection{Open gaps and the data used}
Three gaps remain open, and each can be shown from the work above. First, no dependence rule proposed for
quantum slicing has been checked against what a quantum program computes, and the one example worked out in
detail, the QPDG's, gives a wrong slice. Second, the analyses that represent quantum programs, from code
property graphs to execution-order tables, read the control of a gate as unchanged, and no work has measured
what that reading does to a slice. Third, the slicing support developers have is manual and dynamic, and no
static slicer shows a qubit's slice on Qiskit source with its loop iterations.

\begin{table*}[t]
  \centering
  \caption{Related approaches by what they analyze, whether they slice from a criterion, what they guarantee,
  whether they report source lines, and the data they were evaluated on.}
  \label{tab:related}
  \footnotesize\setlength{\tabcolsep}{3pt}
  \begin{tabular}{@{}>{\raggedright\arraybackslash}p{2.6cm} >{\raggedright\arraybackslash}p{2.4cm}
                  >{\raggedright\arraybackslash}p{2.6cm} >{\raggedright\arraybackslash}p{2.6cm} l
                  >{\raggedright\arraybackslash}p{4.4cm}@{}}
    \toprule
    Approach & Analyzes & Slicing criterion & Stated guarantee & Lines & Evaluated on \\
    \midrule
    srcSlice~\cite{alomari2014srcslice} & classical source, srcML & variable, automatic & none & yes & the Linux kernel, against CodeSurfer \\
    Behler et al., QStatic~\cite{behler2024supporting, behler2025static} & OpenQASM~3, srcML & none & none & yes & the 16 OpenQASM~3 example files \\ % TODO verify: QStatic's own evaluation set
    QPDG~\cite{yu2025qpdg} & quantum programs & variable or qubit & none & -- & examples, not implemented \\
    Kaul et al.~\cite{kaul2023uniform} & Qiskit, OpenQASM & none, queries & none & yes & example programs \\
    \texttt{LightCone}~\cite{qiskit2025} & circuits & observable & keeps the observable's value & no & -- \\
    Metwalli, Van Meter~\cite{metwalli2022debugging, metwalli2024testing} & Qiskit, at run time & none, manual breakpoints & none & yes, dynamic & small example circuits \\ % TODO verify: count of circuits and the survey
    LintQ~\cite{paltenghi2024lintq} & Qiskit, CodeQL & none, bug patterns & none & yes & 7,568 GitHub programs, \numLqTp{} confirmed bugs \\
    QChecker~\cite{zhao2023qchecker} & Qiskit, syntax tree & none, bug patterns & none & yes & the 42 bugs of Bugs4Q \\
    \midrule
    This paper, \tool{} & Qiskit, OpenQASM~2, srcML and srcSlice & a set of qubits, automatic & sound, dependence-complete & yes, static & \numNCircuits{} benchmark circuits, \numFamCircuits{} generated circuits, Bugs4Q, the Qiskit textbook, LintQ's \numScaleSOneFiles{} programs and warnings, QASMBench, Qiskit HumanEval \\
    \bottomrule
  \end{tabular}
\end{table*}

The data differ as much as the methods. The dependence graphs were shown on examples, the srcML work on the
OpenQASM~3 examples, and the detectors on Bugs4Q and on LintQ's GitHub corpus, while no slicing approach was
evaluated on real programs. We use the detectors' data where they fit, Bugs4Q for real faults and LintQ's corpus
and confirmed warnings for real programs, and add what a slicing study needs that they lack, circuits with a
simulated ground truth. These are \numNCircuits{} circuits from MQT Bench~\cite{quetschlich2023mqt}, QASMBench~\cite{li2023qasmbench} and
VeriQBench~\cite{chen2022veriqbench}, \numFamCircuits{} generated circuits with known structure, and seeded faults. For the scale of real programs we
add QASMBench's large circuits and, for current Qiskit, the expert-written programs of Qiskit
HumanEval~\cite{vishwakarma2024qhe}.


\section{Quantum Dependence and the Write-All Rule}\label{sec:core}

This section defines what a slice of a set of qubits must preserve and when an operation is a dependence of
those qubits. It then states the write-all rule, proves that it gives slices that are sound and contain every
dependence, shows that three kinds of dependence are all a slice needs, and places every other rule relative to
it. The definitions do not depend on any slicer, and Section~\ref{sec:results} measures each rule against them.

\subsection{Circuits, criteria and dependence}
A Qiskit program does not act on qubits directly. It builds a \emph{circuit}, which a simulator or a quantum
computer then runs. The circuit $P = o_0 \dots o_{L-1}$ is the sequence of operations the program appends, in the
order it appends them, and position $k$ holds the operation $o_k$. The circuit has $n$ qubits, numbered
$Q = \{0, \dots, n-1\}$, which form its quantum memory, and a set $C$ of classical bits, which store the outcomes of
measurements.

Each operation touches some qubits and bits, which we call its \emph{locations} $\mathit{loc}(o) \subseteq Q \cup
C$. A gate touches the qubits it is applied to, so \texttt{h(0)} has the location $q_0$ and \texttt{cx(0, 1)} the
locations $q_0$ and $q_1$. A measurement touches the measured qubit and the bit that receives its outcome, so
\texttt{measure(2, 0)} has the locations $q_2$ and $c_0$. A reset, which returns a qubit to $\ket{0}$, touches that
qubit. A conditioned gate is applied only when a measured bit has a given value, as in \texttt{x(0)} inside
\texttt{if\_test((c[0], 1))}, and it touches its qubits and the bits its condition reads, here $q_0$ and $c_0$.

The program's own loops and branches run while the circuit is built and decide which operations it contains and
how often. Loops that run inside the circuit, on the device, are treated in Section~\ref{sec:loops}.

The state of the qubits is described by a density matrix $\rho$, which also covers the mixtures that arise when a
measurement outcome is not known in advance. Running $P$ on the input $\rho$, with every bit initially $0$, gives
$\llbracket P \rrbracket(\rho)$, the final joint state of qubits and bits, averaged over the outcomes of the
measurements. A developer asks about some of the qubits, a set $W \subseteq Q$. Their \emph{reduced state}
$\rho_W = \operatorname{Tr}_{\overline{W}}(\rho)$ is what remains of $\rho$ when everything else,
$\overline{W} = (Q \setminus W) \cup C$, is traced out, that is, averaged away. It determines every statistic of
measuring $W$ in any basis, so two runs that give $W$ the same reduced state cannot be told apart by looking at
$W$. A slice is a set $S$ of positions in $P$, and it is judged by running the restriction $P|_S$, the circuit that
keeps only the operations at the positions in $S$, in their original order.

\begin{definition}[Slicing criterion]\label{def:criterion}
A slicing criterion is a nonempty set $W \subseteq Q$ of qubits, whose state is taken at the end of the circuit,
before any terminal measurement reads it.
\end{definition}

The criterion is a set of qubits, not a single one, because a qubit has no value of its own apart from the qubits
it is correlated with. Yu and Zhao keep the classical criterion of a variable at a statement~\cite{yu2025qpdg}, and
we preserve the \emph{joint} state of the criterion instead. Taking that state before readout keeps a slice correct
when it is composed with later gates or read in another basis.

\begin{definition}[Soundness]\label{def:sound}
A set of positions $S$ \emph{preserves} $W$ at input $\rho$ if $\operatorname{Tr}_{\overline{W}}
\llbracket P|_S \rrbracket(\rho) = \operatorname{Tr}_{\overline{W}}\llbracket P \rrbracket(\rho)$. It is a
\emph{sound} slice for $W$ if it preserves $W$ at every input.
\end{definition}

A sound slice reproduces the criterion, but comprehension needs more, the assurance that the slice shows every
place where a change, and so a fault, could alter the criterion.

\begin{definition}[Dependence]\label{def:dep}
Operation $o_k$ is a \emph{dependence} of $W$ in $P$ if there are an input $\rho$ and a unitary $U$ on the qubits
of $o_k$ such that inserting $U$ immediately after $o_k$ changes $\operatorname{Tr}_{\overline{W}}\llbracket P
\rrbracket(\rho)$. For a measurement, negating the bit it writes also counts as such a change. $\mathit{Dep}_W(P)$
is the set of dependences of $W$, and a set $S$ is \emph{dependence-complete} for $W$ if $\mathit{Dep}_W(P)
\subseteq S$.
\end{definition}

This is semantic dependence in the sense of Podgurski and Clarke~\cite{podgurski1990formal}, stated for quantum
states. The two guarantees differ. A sound slice need not contain every dependence. An identity gate on a qubit
of $W$ is a dependence, since a perturbation inserted after it changes the criterion, yet deleting it changes
nothing. Conversely, deleting one unitary operation that is not a dependence never changes $\rho_W$, as
inserting its inverse shows.

\subsection{The write-all rule}
A \emph{dependence rule} $R$ assigns each operation the locations it reads, $r_R(o)$, and the locations it writes,
$w_R(o)$. Operation $o_j$ depends on $o_i$ under $R$ if $i < j$ and some location in $w_R(o_i) \cap r_R(o_j)$ is
not written by an operation between them. The backward slice $B_R(W)$ is the set of operations from which the
last writer of some $q \in W$ is reachable along these dependences. The rules differ in one choice, what an
operation writes.

\begin{definition}[Write-all rule]\label{def:writeall}
Every operation reads and writes all of its locations, $r(o) = w(o) = \mathit{loc}(o)$. In particular a
controlled gate writes its controls as well as its targets. The \emph{quantum dependence graph} (QDG) has the
operations as nodes and the dependences of this rule as edges, and the QDG slice $\Bw(W)$ is backward
reachability from $W$.
\end{definition}

The QDG slice is computed in one backward pass. Start at the end of the circuit with the needed set $N = W$, and
whenever an operation acts on a location in $N$, keep the operation and add all its locations to $N$. On qubits
alone, this is the backward light cone of $W$. The classical bits extend it, from a measurement to the
operations its bit conditions. Take the circuit
\[
P_1 \;=\; \mathrm{X}(q_1)\;\; \mathrm{H}(q_1)\;\; \mathrm{H}(q_0)\;\; \mathrm{CX}(q_0, q_1)
\]
and the criterion $\{q_0\}$. The CNOT acts on $q_0$, so it is kept and $q_1$ joins $N$, which keeps the two gates
on $q_1$ as well. Although $q_0$ is the CNOT's control, the gates on $q_1$ do change it. At input $\ket{00}$ the
target is prepared in $\ket{-}$, an eigenstate of the bit flip, so the CNOT leaves the target as it is and flips
the phase of the control, which ends in $\ket{-}$ instead of $\ket{+}$. This is phase kickback. A rule that
reads the control as unchanged lets the CNOT write only $q_1$, so the last operation to write $q_0$ is
$\mathrm{H}(q_0)$ and its slice keeps that gate alone.

\begin{theorem}[The QDG is sound]\label{thm:sound}
For every circuit $P$ and criterion $W$, $\Bw(W)$ is a sound slice for $W$.
\end{theorem}

\begin{IEEEproof}[Proof sketch]
Process $P$ backward with the needed set $N$, initially $W$. Operation $o_k$ joins the slice, and $N$ grows by
$\mathit{loc}(o_k)$, exactly when $\mathit{loc}(o_k) \cap N \neq \emptyset$, which computes $\Bw(W)$. Every gate,
reset, measurement and conditioned gate is a channel that acts only on $\mathit{loc}(o)$. By locality of the partial
trace, the no-signaling principle, an operation outside the slice cannot change the reduced state of the needed
locations after it, and an operation in the slice maps the reduced state of the needed set before it to the one
after it. Induction over $P$ gives the claim. The induction is mechanized in Isabelle/HOL for an abstract model of
local operations, and that quantum operations satisfy the model is argued on paper.
\end{IEEEproof}

\begin{proposition}[The QDG is dependence-complete]\label{prop:complete}
$\mathit{Dep}_W(P) \subseteq \Bw(W)$.
\end{proposition}

\begin{IEEEproof}
Let $o_k \notin \Bw(W)$, and let $P'$ insert a unitary $U$ after $o_k$. In the backward pass, no location of
$o_k$ is needed after it, since otherwise $o_k$ would be kept. $U$ acts only on locations of $o_k$, so $U$ is not
kept either and the needed set is unchanged. Hence $P'$ and $P$ have the same QDG slice, and by
Theorem~\ref{thm:sound} the same reduced state of $W$. A negated outcome is handled alike, as an operation on the
bit of a measurement outside the slice.
\end{IEEEproof}

\begin{proposition}[The QDG distributes over criteria]\label{prop:joint}
$\Bw(W_1 \cup W_2) = \Bw(W_1) \cup \Bw(W_2)$.
\end{proposition}

\begin{IEEEproof}
In the backward pass of Theorem~\ref{thm:sound}, the needed set for $W_1 \cup W_2$ is at every position the union
of those for $W_1$ and $W_2$.
\end{IEEEproof}

Together these say that the QDG slice is exactly what a slice must keep. It keeps everything that can change the
criterion and reproduces the criterion for every input, and the slice of a set of qubits is the union of the slices
of its members, so a tool computes single-qubit slices and combines them. The relation itself is not new. It is the
light cone that compilers compute, and LintQ orders operations by it~\cite{paltenghi2024lintq}. What the three
results add is the justification of that relation as a slice.

\paragraph*{Why the write set cannot be smaller}
Dropping a qubit from the write set of a two-qubit gate is safe only if that gate can never carry information from
the other qubit to it. No standard gate qualifies. For each of the \numGateSignalN{} two-qubit gates of Qiskit's
standard library, from the CNOT and CZ to the SWAP and the parameterized rotations, a change of either qubit's input
changes the reduced state of the other for some product input (exploratory, Deviation I-24). A rule that reads any
of these gates as writing only some of its qubits therefore misses a dependence of some circuit.
% TODO authors: a general necessity proposition (every two-qubit gate that is not a product of single-qubit gates
% signals in both directions) needs checking against Beckman et al., Phys. Rev. A 64 (2001), before it is stated.

\subsection{Three kinds of dependence suffice}\label{sec:kinds}
The write-all rule needs only three kinds of dependence, which Table~\ref{tab:kinds} lists with examples from the
ansatz of Fig.~\ref{fig:ansatz}. \emph{Classical data} dependence follows the host values that compute gate
arguments, such as angles, sizes and qubit indices. \emph{Classical control} dependence follows what decides
whether an operation is applied at all. It covers the host program's loops and branches, which decide when the
circuit is built which operations it contains and how often, and the measured bits that condition operations when
the circuit runs, in a conditioned gate, a branch or a loop of the circuit. \emph{Quantum data} dependence follows
a qubit through the operations that act on it, under the write-all rule.

A controlled gate looks like a fourth kind, quantum control dependence, since it reads like a branch, if the control
is 1 then flip the target. It is not one. Its control is coherent, so the gate acts on a superposition of control
values and no branch is ever taken, and it can change its control, as $P_1$ shows. What looks like control
dependence is quantum data dependence in both directions, which the write-all rule records, and reading the gate as
a branch is exactly the error of the classical PDG in Section~\ref{sec:rules}. Real control dependence in a quantum
program arises only through measurement, whose classical bit then conditions later operations.

\begin{table}[t]
  \centering
  \caption{The kinds of dependence a quantum slice needs}
  \label{tab:kinds}
  \footnotesize\setlength{\tabcolsep}{3pt}\renewcommand{\arraystretch}{1.15}
  \begin{tabular}{@{}>{\raggedright\arraybackslash}p{1.55cm} >{\raggedright\arraybackslash}p{3.2cm}
                  >{\raggedright\arraybackslash}p{3.15cm}@{}}
    \toprule
    \textbf{Kind} & \textbf{Follows} & \textbf{In Fig.~\ref{fig:ansatz}a} \\
    \midrule
    Classical data & host values that compute gate arguments & \texttt{n}, \texttt{layers}, the angles \\
    \arrayrulecolor{black!25}\hline\arrayrulecolor{black}
    Classical control & host loops and branches, and measured bits that condition operations & \texttt{for qubit in range(n)}, the flip conditioned on the measured bit \\
    \arrayrulecolor{black!25}\hline\arrayrulecolor{black}
    Quantum data & every qubit an operation acts on, written by the write-all rule & \texttt{cx(qubit, qubit+1)} writes both qubits \\
    \midrule
    Quantum control & no kind of its own, it is quantum data dependence & a controlled gate is not a branch \\
    \bottomrule
  \end{tabular}
\end{table}

The other kinds proposed for quantum programs add no dependence either. Entanglement is the most natural
candidate, since entangled qubits are correlated, but it decides nothing.

\begin{proposition}[Entanglement is neither necessary nor sufficient for dependence]\label{prop:ent}
(i) An operation can be a dependence of $q$ although, at an input that witnesses the dependence, $q$ is never
entangled with its qubits. (ii) An operation can act on a qubit entangled with $q$ and not be a dependence of $q$.
\end{proposition}

\begin{IEEEproof}
(i) In $P_1$ at input $\ket{00}$, the state is a product state after every operation, yet the gates on $q_1$ are
dependences of $q_0$, since inserting $\mathrm{H}(q_1)$ after them changes the final state of $q_0$. (ii) In
$\mathrm{H}(q_0)\,\mathrm{CX}(q_0, q_1)\,\mathrm{X}(q_1)$ the two qubits are maximally entangled before
$\mathrm{X}(q_1)$, but by no-signaling no operation on $q_1$ after the CNOT can change $\rho_{q_0}$.
\end{IEEEproof}

By Proposition~\ref{prop:complete}, entanglement edges can add no dependence that the QDG misses, and by (ii) they
can add operations that are not dependences. The remaining proposed kinds add no edges. A CNOT that copies a basis
state, the case no-cloning concerns, is a two-qubit gate whose edges quantum data dependence already has. Timing
decides when an operation runs on hardware, not what a state depends on, and a compiler can move a single-qubit
gate within the idle time of its qubit without changing the ideal computation~\cite{smith2022timestitch}.
Superposition changes amplitudes, not which locations an operation acts on~\cite{nielsen2010quantum}, so it cannot
add an edge to a dependence graph.

\subsection{The other rules, relative to the write-all rule}\label{sec:rules}
The study compares the write-all rule with the rules a developer or tool builder would otherwise use.
Table~\ref{tab:rules} gives each by how it differs from the QDG, and Fig.~\ref{fig:ansatz} applies them to one
program that needs every kind of dependence. An ansatz is the parameterized circuit of a variational algorithm,
whose rotation angles a classical optimizer tunes. The ansatz of Fig.~\ref{fig:ansatz}a has layers of rotations and
a chain of CNOTs, and after them an auxiliary qubit is measured mid-circuit and $q_0$ is flipped when the outcome
is~$1$.

\begin{figure*}[t]
  \centering
  \includegraphics[width=\textwidth]{figures/fig_pdg_ansatz.pdf}
  \caption{The rules on a variational ansatz, criterion $q_0$. (a)~The program. After the loops, an auxiliary qubit
  is measured mid-circuit and conditions a flip of $q_0$. (b)~One graph for the first four rules of
  Table~\ref{tab:rules}. Each node is an operation as it executes, with its source line and loop iteration, and a box
  holds each iteration of the outer loop. Grey data edges belong to every rule, and def-use lacks the green control
  edge from the measurement. The classical PDG and def-use lack the blue quantum edges and have the dotted edges
  instead, each of which skips a CNOT whose control it reads as unchanged. A node is shaded by the weakest rule that
  keeps it. The QDG keeps exactly the dependences of $q_0$, and the QPDG adds the hatched operations, which are not
  dependences. Graphs, slices and fidelities at $\ket{0\dots0}$ are computed by the study's code from the program
  shown.}
  \label{fig:ansatz}
\end{figure*}

\begin{table}[t]
  \centering
  \caption{The dependence rules, by what a CNOT with control $c$ and target $t$ writes and by how each differs from
  the QDG}
  \label{tab:rules}
  \footnotesize\setlength{\tabcolsep}{2.6pt}\renewcommand{\arraystretch}{1.12}
  \begin{tabular}{@{}l c >{\raggedright\arraybackslash}p{2.05cm} >{\raggedright\arraybackslash}p{1.38cm} r r@{}}
    \toprule
    \textbf{Rule} & $w(\mathit{cx}_{c,t})$ & \textbf{Relation to the QDG} & \textbf{Guarantee} & \textbf{Ops} & $F$ \\
    \midrule
    \rowcolor{blue!7}\textbf{QDG} & $\{c, t\}$ & \textbf{write-all rule} & \textbf{sound, complete} & \numAnsWireOps{} & \numAnsQdgFidelity{} \\
    \midrule
    Def-use & $\{t\}$ & $-$ control write, $-$ bits & none & \numAnsDuOps{} & \numAnsDuFidelity{} \\
    Classical PDG & $\{t\}$ & $-$ control write & none & \numAnsCpdgOps{} & \numAnsCpdgFidelity{} \\
    QPDG~\cite{yu2025qpdg} & $\{c, t\}$ & $+$ closed entanglement edges & sound & \numAnsYzOps{} & \numAnsYzFidelity{} \\
    QPDG, read-only & $\{t\}$ & $-$ control write, $+$ entanglement & none & \numAnsYzroOps{} & \numAnsYzroFidelity{} \\
    \bottomrule
  \end{tabular}\\[3pt]
  \begin{minipage}{\columnwidth}\footnotesize
  Every rule reads $\{c, t\}$, and complete means dependence-complete. Commutation and state-based pruning remove
  operations from the QDG slice. Ops is the number of the \numAnsOps{} operations of Fig.~\ref{fig:ansatz} that the
  slice of $q_0$ keeps, and $F$ its fidelity at $\ket{0\dots0}$.
  \end{minipage}
\end{table}

\paragraph*{The classical PDG and def-use}
The classical PDG, written CPDG in subscripts, applies the program dependence graph~\cite{ferrante1987pdg} to the
circuit. It reads a gate as a classical assignment, so a controlled gate reads its controls and targets and writes
only its targets, and any other gate reads and writes all its qubits. A measurement writes its bit, and a
conditioned gate reads the bits its condition tests. \emph{Def-use} (DU) is the same data dependence without the
bits, the simplest transfer of classical slicing. Reading a control as unchanged is the usual explanation of a
CNOT~\cite{ramalho2025road}, and it holds for computational basis states. A table of the operations on each qubit
that marks the control of a CNOT apart from its target~\cite{behler2024supporting} serves pattern detection well,
but read as a dependence graph it gives this reading.

\begin{proposition}[The classical PDG is unsound]\label{prop:du}
There are circuits for which neither $B_{\mathrm{CPDG}}(W)$ nor $B_{\mathrm{DU}}(W)$ preserves $W$, even at
$\ket{0\dots0}$.
\end{proposition}

The ansatz of Fig.~\ref{fig:ansatz} is a witness for both, and on unitary circuits the two slices are equal.
Def-use ignores the classical bits, so its slice of $q_0$ keeps only \numAnsDuOps{} of the \numAnsOps{} operations,
with fidelity \numAnsDuFidelity{}, and loses the auxiliary's gate and measurement, which the classical PDG reaches
through the green control edge in Fig.~\ref{fig:ansatz}b. Control dependence is not enough, however. Each
\texttt{cx(qubit, qubit+1)} changes its control as well as its target, so the state of $q_0$ comes to depend on the
rotations of $q_1$ and, from the second layer on, of $q_2$. The classical PDG keeps only \numAnsCpdgOps{}
operations, misses \numAnsCpdgMissed{} dependences and reaches fidelity \numAnsCpdgFidelity{}. Where the QDG has
a blue quantum edge out of a CNOT, the classical PDG has a dotted edge that skips the CNOT. The QDG slice keeps the
\numAnsWireOps{} operations that are dependences of $q_0$, and no other.

\paragraph*{The QPDG}
Our reimplementation of Yu and Zhao's graph~\cite{yu2025qpdg} serves as a baseline. Its dependences are given in
prose, so we logged every reading of it before any run. Data and control dependence are as in the QDG. The QPDG
gives no rule for when two qubits are entangled, so we use the flow-insensitive rule of entanglement analysis, under
which qubits joined by a two-qubit gate anywhere in the circuit are entangled~\cite{javadiabhari2015scaffcc}, and
add an edge from every operation to every later one on qubits of the same component. A flow-sensitive variant
counts two qubits as entangled only from the first gate that couples them. Either way the QPDG is the QDG plus
entanglement edges, so its slices contain those of the QDG, and in Fig.~\ref{fig:ansatz}b the extra edges add the
\numAnsYzExtra{} hatched operations of $q_2$ in the second layer, which run after the last gate that carries
information from $q_2$ to $q_0$. Yu and Zhao do not say whether a gate writes its control, so we also run, as an
exploratory variant, the reading in which it does not.

\paragraph*{Pruning}
Two refinements start from the QDG slice and remove operations from it. \emph{Commutation pruning} propagates the
Pauli observables $X$, $Y$ and $Z$ of every criterion qubit backward and removes each operation that commutes with
all of them, since such an operation cannot change a measurement statistic of the criterion. Qiskit's
\texttt{LightCone} pass implements it~\cite{qiskit2025}. \emph{State-based pruning} fixes one input, by default
$\ket{0\dots0}$, and removes an operation that does not act on the criterion if the mutual information between the
criterion and the operation's qubits is below $10^{-9}$ and the removal still preserves $W$ at that input. Both give
smaller slices, but neither is dependence-complete, and a state-based slice holds for its one input only, as a
conditioned slice does~\cite{canfora1998conditioned}.

\subsection{Loops}\label{sec:loops}
A loop of the host program runs when the circuit is built. A static slicer unrolls it, so every iteration
contributes its own operations, each tagged with the line and iteration that emit it. When the number of iterations
cannot be known statically, the loop body is repeated as \emph{may-operations}, operations that might be emitted,
until no repetition adds a location to any dependence chain. A loop of the circuit itself, Qiskit's
\texttt{while\_loop}, \texttt{for\_loop} or \texttt{switch}, runs on the device, and its condition reads measured
bits. Its body is treated the same way, with every operation in it conditioned on the bits its condition reads, so
the classical control dependence on the measurement that decides the loop is kept.

\begin{proposition}[Loops]\label{prop:loops}
Let a loop body $B$ act on a set of locations $\Lambda$, let the loop end with probability one, and let $P_B$ repeat
$B$ $|\Lambda| + 1$ times, with every operation of $B$ conditioned on the bits that decide the loop. Deleting from
$B$ every operation that no copy in the QDG slice of $P_B$ keeps gives a sound slice for every criterion.
\end{proposition}

\begin{IEEEproof}[Proof sketch]
A loop that runs at most $K$ times is the circuit of $K$ copies of $B$, each conditioned on the bits that decide
the loop, a circuit with conditioned operations to which Theorem~\ref{thm:sound} applies. In its backward pass every
copy is processed alike, and the needed set only grows. Once a copy adds no location of $\Lambda$, every earlier copy
keeps the same operations, and the needed set can grow at most $|\Lambda|$ times, so for every $K$ the copies keep
only operations that one of the $|\Lambda| + 1$ copies of $P_B$ keeps. The bits that decide the loop are read by
every copy, so the measurements that write them are kept and the distribution over numbers of iterations is
preserved. A loop that ends with probability one is the limit of its bounded prefixes.
\end{IEEEproof}

The induction of Theorem~\ref{thm:sound} is not mechanized for loops, and the proposition is argued on paper. It is
checked exactly on \numCfPrograms{} programs of our own with circuit loops, branches and switches. For every one of
their \numCfQubits{} qubits, deleting every operation outside the qubit's slice from the source leaves the qubit's
state unchanged at \numCfInputs{} inputs, while the slices that read a circuit loop's body once are wrong on
\numCfOldWrong{} of the \numCfPrograms{} programs (exploratory, Deviation I-26).

\subsection{Forward slices and qubit profiles}
The forward slice $\Fw(o_k)$ of an operation is the set of operations reachable from $o_k$ in the QDG. It answers
the question a developer asks before changing a gate, which later operations and qubits the change can reach. We say
that $o_k$ reaches $q$ if $\Fw(o_k)$ contains the last operation on $q$.

\begin{proposition}[Forward and backward slices are dual]\label{prop:dual}
For every operation $o_k$ and qubit $q$, $o_k$ reaches $q$ if and only if $o_k \in \Bw(\{q\})$. Hence $o_k$ reaches
every qubit of which it is a dependence.
\end{proposition}

\begin{IEEEproof}
Both sides say that the QDG has a path from $o_k$ to the last operation on $q$. The second claim follows from
Proposition~\ref{prop:complete}.
\end{IEEEproof}

A slice can be reported as operations, loop iterations or source lines. Operations are the unit of the definitions
and lines the unit of the developer. A line is in a slice if one of its operations is, so a line can be kept for
one iteration only. In the second layer of the ansatz, the rotations of $q_2$ are not dependences of $q_0$, yet
their lines are in its slice because their first-layer iterations are. A line slice therefore names the iterations
that put each line in it.

The slice profile of a variable, which srcSlice records~\cite{alomari2014srcslice}, has a quantum counterpart that
gathers in one view what a developer first asks about a qubit.

\begin{definition}[Qubit slice profile]\label{def:profile}
The profile of a qubit $q$ records its \emph{uses}, the lines of the operations on $q$, each of which both reads
and writes $q$ under the write-all rule; its \emph{names}, the source expressions that denote $q$ in gate calls; its
\emph{functions}, the call sites of the functions in which $q$ is acted on; its \emph{dependents}, the qubits that the
first operation on $q$ reaches; its \emph{bits}, those measured from $q$, with the lines they condition; and its
\emph{backward} and \emph{forward} entries, the lines of the backward slice of $q$ and of the forward slice of its
first operation.
\end{definition}


\section{Study Design}\label{sec:design}

The study measures every rule of Section~\ref{sec:core} against its definitions, asking four questions a developer
has of a slice.

\begin{itemize}
\item \textbf{RQ1} Which dependence rules give sound and dependence-complete slices?
\item \textbf{RQ2} How much smaller than the QDG slice can a slice safely be, and under which inputs?
\item \textbf{RQ3} What does each rule do to fault localization?
\item \textbf{RQ4} What does each rule cost in time as circuits grow?
\end{itemize}

\subsection{Pre-registration}
The protocol was committed before the runs reported here. It extends an earlier run whose results it reuses, so the
outcomes of H1 and H2 were known when it was written, and every later change is a dated deviation. The five
confirmatory tests concern state-based pruning. Every other result is exploratory or descriptive, including the
classical PDG's share of wrong slices and everything added after the first results, the classical PDG and
\texttt{LightCone} baselines, the dependence set and its nesting with the rules, the symbolic-angle, flow-sensitive and
read-only variants, and the gate check. Two descriptive tasks for language models were pre-registered, and the most
expensive model was dropped for cost before any scored call.

\subsection{Subjects}
The rules are compared on circuits, where exact simulation gives the ground truth. Corpus~A collects
\numSubjACount{} benchmark circuits from MQT Bench~\cite{quetschlich2023mqt}, QASMBench~\cite{li2023qasmbench} and
VeriQBench~\cite{chen2022veriqbench}, with \numSubjAMinQubits{} to \numSubjAMaxQubits{} qubits and a median of
\numSubjAMedianOps{} operations, at most \numSubjAMaxOps{}. Every qubit of every circuit is a criterion, which gives
\numNPairs{} circuit--qubit pairs, \numNDynamicPairs{} of them in dynamic circuits with mid-circuit measurement, reset
or conditioned gates. The \numSubjGenCount{} generated circuits come in families of known structure that grow in depth
and in width, with \numSubjGenMinQubits{} to \numSubjGenMaxQubits{} qubits and up to \numSubjGenMaxOps{} operations,
and \numSubjSweepCount{} width sweeps extend two of the families to \numSubjSweepMaxQubits{} qubits. Corpus~L has
\numSubjLCount{} wider circuits from the same suites, with \numSubjLMinQubits{} to \numSubjLMaxQubits{} qubits and up
to \numSubjLMaxOps{} operations, on which only some rules finish. Faults are seeded into corpus~A, and real faults
come from Bugs4Q~\cite{zhao2023bugs4q}. A unit is a fault and a qubit whose measured distribution it changes.

\subsection{Ground truth}
An independent oracle simulates each circuit exactly, branching over measurement outcomes, and agrees with Qiskit's
simulation to within \numOracleErr{}. A slice preserves the criterion when the squared Uhlmann fidelity of the
criterion's reduced states under the slice and the whole circuit is at least $1 - 10^{-6}$. The
dependences of Definition~\ref{def:dep} are measured by inserting \numInfluencePerturbations{} random unitaries after
every operation and comparing the criterion's state at \numInfluenceInputs{} inputs, $\ket{0\dots0}$ and random
product states. Being sampled, the measured dependence set can only miss true dependences. A catalog of \numCatEntries{} hand-derived cases checks the ground truth and every
slicer before the runs.

\subsection{Slicers}
Each rule is implemented on circuits from its read and write sets or its pruning step. Qiskit's \texttt{LightCone}
pass applies to \numLcqApplicable{} of the \numNPairs{} pairs. It rejects mid-circuit measurements on
\numLcqNaMeasurements{} pairs, and on \numLcqNaReorder{} its output cannot be matched back to the original.
Exploratory variants make commutation pruning hold for all angles and give state-based pruning a joint criterion.

\subsection{Robustness}
A slice computed at $\ket{0\dots0}$ is rechecked, with its operations unchanged, under \numRobDraws{} random
computational-basis inputs, \numRobDraws{} random product inputs and, for circuits with rotation angles,
\numRobDraws{} random rebindings of every angle. A pair is robust if the slice preserves the criterion in every draw. The gate check takes each of the
\numGateSignalN{} two-qubit gates of Qiskit's standard library, parameterized gates at a fixed angle, and searches
seeded random product inputs, for each direction, for two inputs of one qubit that leave the other in different
reduced states, evidence for these gates and not a proof for all.

\subsection{Faults and fault localization}
Faults are seeded with \numMutOperators{} operators, \numMutPublished{} from Muskit~\cite{mendiluze2021muskit},
QMutPy~\cite{fortunato2022qmutpy} and García de la Barrera et al.~\cite{garcia2022testing}, and an angle perturbation
of our own. Of the \numMutants{} mutants, \numMutantsExcluded{} with an empty fault site are excluded, and
\numMutantsKilled{} change the measured distribution of at least one qubit. Each such qubit and its mutant form a
unit, \numRqThreeUnits{} in all, for which we record each slice's size and whether it contains the fault site. Of the
\numBugsFourQTriaged{} entries of Bugs4Q, \numBugsFourQIncluded{} are circuit-level faults that port mechanically to
the current Qiskit release, and \numBugsFourQKilled{} of them change a qubit's measured distribution, giving
\numLinesUnits{} units. Running each port
records the line that appends each operation, mapped to the original program, so these slices are also reported as
lines of the programs their developers wrote.

Mutation-based fault localization ranks the operations of the faulty circuit by how their mutants change its
observable behavior, with the Ochiai formula, as Ishimoto et al.\ do for quantum programs~\cite{ishimoto2025mbfl}, on a
sample of \numMbflUnits{} units stratified by fault operator and including every Bugs4Q unit. We rank over the whole
circuit and within each slice and report the EXAM score, the share of ranked operations examined before the fault,
lower being better.

In the language-model tasks, a model lists the operations that can affect a qubit for \numTOneItems{} items, and
names the likely faulty operation of \numTTwoItems{} units given the whole circuit or a slice. The models are Claude
Sonnet~5.5 and Haiku~5.5 and Qwen2.5-Coder with 32 and 1.5 billion parameters.

\subsection{Cost}
RQ4 times the rules on the generated family that grows in depth at eight qubits, on the width sweeps and on
corpus~L, each circuit in a fresh process limited to \numCostWallLimit{}~s and 4~GiB.

\subsection{Statistics}
Five confirmatory hypotheses concern state-based pruning. It is wrong on fewer pairs than def-use (H1), its slices
are smaller than QDG slices (H2), at least 95\% of pairs keep a correct slice under random product inputs (H3), it
loses at most five points of fault containment against the QDG slice (H4), and its slices are smaller where both
contain the fault (H5). It starts from the QDG slice and checks every removal at $\ket{0\dots0}$, so H1 and the
directions of H2 and H5 hold by construction, and H3 and H4 are the real tests. We use exact McNemar tests on paired outcomes and Wilcoxon signed-rank tests on paired
sizes, with Holm's correction across the five. Effect sizes come with bootstrap 95\% intervals and Wilson intervals
for proportions, and exploratory intervals resample whole algorithm families. Everything else is descriptive.


\section{Results}\label{sec:results}

Table~\ref{tab:scorecard} gives the answer first. The QDG does what Section~\ref{sec:core} predicts on every pair the
oracle can check, comes within a few points of the smallest complete slice, holds under noise, and costs milliseconds.
It is not guaranteed to contain a fault, since a fault can delete the very interaction a slice would follow, and RQ3
measures how often that happens. The subsections answer each question for every rule.

\begin{table}[t]
  \centering
  \caption{The QDG on every question of the study}
  \label{tab:scorecard}
  \footnotesize\setlength{\tabcolsep}{3pt}\renewcommand{\arraystretch}{1.12}
  \newcommand{\sep}{\arrayrulecolor{black!25}\hline\arrayrulecolor{black}}
  \begin{tabular}{@{}l >{\raggedright\arraybackslash}p{2.0cm} >{\raggedright\arraybackslash}p{4.15cm}@{}}
    \toprule
    \textbf{Question} & \textbf{Expected} & \textbf{Observed} \\
    \midrule
    Sound (RQ1) & every input, Theorem~\ref{thm:sound} & \numRqOneLcUnsound{} wrong slices of \numNPairs{} pairs, and of \numCorpusLSoundPairs{} on corpus~L \\ \sep
    Complete (RQ1) & Proposition~\ref{prop:complete} & recall \numRqOneLcRecall{} of the measured dependences \\ \sep
    Size (RQ1) & not minimal & \numLcExcess{} points above the dependence set, equal to it on \numLcExact{} pairs \\ \sep
    Noise (RQ2) & local noise, Theorem~\ref{thm:sound} & trace distance at most \numNoiseLcCrossExactMax{} \\ \sep
    Faults (RQ3) & no guarantee & contains \numRqThreeContainLc{} of the faults, \numSizeCtlLcGain{} points above a random set of its size \\ \sep
    Cost (RQ4) & one pass & at most \numCostLcMaxCritMs{}~ms per criterion \\
    \bottomrule
  \end{tabular}
\end{table}

\subsection{RQ1: Which rules give sound and complete slices?}

\begin{table}[t]
  \centering
  \caption{Slices of the \numNPairs{} corpus-A pairs for RQ1}
  \label{tab:rq1}
  \footnotesize\setlength{\tabcolsep}{3pt}\renewcommand{\arraystretch}{1.12}
  \newcommand{\sep}{\arrayrulecolor{black!25}\hline\arrayrulecolor{black}}
  \begin{tabular}{@{}l r r r r r@{}}
    \toprule
    \textbf{Rule} & \textbf{Wrong} & \textbf{Dynamic} & \textbf{Recall} & \textbf{Precision} & \textbf{Size} \\
    \midrule
    Dependence set$^{a}$ & -- & -- & 1.000 & 1.000 & \numDepFloor{} \\
    \midrule
    \rowcolor{blue!7}\textbf{QDG} & \textbf{\numRqOneLcUnsound{}} & \numLcDynamicUnsound{} & \textbf{\numRqOneLcRecall{}} & \numRqOneLcPrecision{} & \numRqOneLcSize{} \\ \sep
    Def-use & \numRqOneDuUnsound{} & \numDuDynamicUnsound{} & \numRqOneDuRecall{} & \numRqOneDuPrecision{} & \numRqOneDuSize{} \\ \sep
    Classical PDG & \numCpdgUnsound{} & \numCpdgDynamicUnsound{} & \numCpdgRecall{} & \numCpdgPrecision{} & \numCpdgSize{} \\ \sep
    QPDG~\cite{yu2025qpdg} & \numRqOneYzUnsound{} & \numYzDynamicUnsound{} & \numRqOneYzRecall{} & \numRqOneYzPrecision{} & \numRqOneYzSize{} \\ \sep
    QPDG, flow-sensitive$^{b}$ & \numYzfUnsound{} & \numYzfDynamicUnsound{} & \numYzfRecall{} & \numYzfPrecision{} & \numYzfSize{} \\
    \midrule
    Commutation$^{c}$ & \numLcqUnsound{} & \numLcqDynamicUnsound{} & \numLcqRecall{} & \numLcqPrecision{} & \numLcqSize{} \\ \sep
    Commutation, symbolic$^{c}$ & \numLcqsUnsound{} & \numLcqsDynamicUnsound{} & \numLcqsRecall{} & \numLcqsPrecision{} & \numLcqsSize{} \\ \sep
    State-based & \numRqOneQUnsound{} & \numQDynamicUnsound{} & \numRqOneQRecall{} & \numRqOneQPrecision{} & \numRqOneQSize{} \\
    \bottomrule
  \end{tabular}\\[3pt]
  \begin{minipage}{\columnwidth}\footnotesize
  Wrong counts the slices that do not preserve the criterion at $\ket{0\dots0}$, a lower bound for every rule but the
  QDG, which is sound for all inputs, and Dynamic those among the \numNDynamicPairs{} pairs in dynamic circuits.
  Recall and precision are against the dependence set, and size is the mean share of the circuit. $^{a}$\,The
  smallest slice that contains every dependence, a reference and not a rule. $^{b}$\,Exploratory, the variant of
  Section~\ref{sec:rules}. $^{c}$\,On the \numLcqApplicable{} pairs where Qiskit's pass
  applies, \numLcqDynamicN{} of them dynamic.
  \end{minipage}
\end{table}

One choice decides soundness, what a controlled gate writes. The classical PDG, which lets a CNOT write only its
target, misses \numCpdgMissTotal{} dependences over the \numNPairs{} pairs, and all but \numCpdgMissFeedback{} of
them, \numCpdgMissControlPct{}, follow from reading a control as unchanged. Of the misses, \numCpdgMissSameWire{} are
controlled gates whose control is the criterion qubit, \numCpdgMissKickback{} reach the criterion only through such
a gate, \numCpdgMissWritten{} reach it through a control further along the chain, and only \numCpdgMissFeedback{}
pass through classical feedback. As a result, Table~\ref{tab:rq1} shows, the classical PDG gives a wrong slice for
\numCpdgUnsound{} of the pairs, or \numCpdgUnsoundPct{}, with a 95\% interval of \numCcCpdgUnsound{} when the
\numCcFamilies{} algorithm families are resampled as clusters. Its control dependence matters only on dynamic
circuits, where it is wrong on \numCpdgDynamicUnsound{} pair against \numDuDynamicUnsound{} for def-use.

Many of these wrong slices look right. Of the \numCpdgUnsound{}, \numFdCpdgZMatch{}, or \numFdCpdgZMatchPct{}, give
the criterion the same distribution of outcomes in the computational basis as the whole circuit, so only
\numFdCpdgObservable{} pairs, \numFdCpdgObservablePct{} of all, are wrong in what a terminal measurement reads. The
others are wrong in phase, which a histogram check accepts but which changes the outcome once the slice is composed
with later gates or measured in another basis. Testing meets the same limit, since phase-flip faults escape
measurement in the computational basis~\cite{oldfield2024faster}. A developer who checks a slice by its outcomes
would therefore accept almost half of the wrong ones. As H1 expects by construction, state-based pruning, which
starts from the QDG slice, is wrong on fewer pairs than def-use, \numRqOneQUnsound{} against \numRqOneDuUnsound{}
(exact McNemar, Holm \numHOnePHolm{}).

The QDG is sound on every pair and recovers every measured dependence, as Theorem~\ref{thm:sound} and
Proposition~\ref{prop:complete} predict. The QPDG is sound too, since its slices contain the QDG's, but its
entanglement edges add only operations that are not dependences, and its slices grow from \numRqOneLcSize{} to
\numRqOneYzSize{} of the circuit. The flow-sensitive variant still keeps \numYzfExcess{} points above the dependence
set, and a graph with one entanglement edge per two-qubit gate gives exactly the QDG slice on all
\numFdFourEdgeNoClosureEqual{} pairs it applies to, so the excess comes from closing these edges transitively. Under the
reading in which a controlled gate writes only its targets, the QPDG is wrong on \numYzRoUnsound{} pairs. On the
\numCorpusLSoundPairs{} pairs of corpus~L that the oracle can check exactly, the QDG and the QPDG are sound everywhere,
while def-use is wrong on \numCorpusLWrongDu{}.

\begin{figure}[t]
  \centering
  \includegraphics[width=\columnwidth]{figures/fig_size.pdf}
  \caption{Mean slice size by circuit width, for RQ1. The dependence set is the floor that every complete slice
  contains. Corpus~L has no dependence set, since measuring one there would need more simulation than the oracle can
  do exactly.}
  \label{fig:size}
\end{figure}

The rules nest, and Fig.~\ref{fig:size} shows how far apart they are. On all \numNestPairs{} pairs the dependence set
lies within the QDG slice and the QDG slice within the QPDG slice. Restricted to its dependence set, every circuit
preserves the criterion at $\ket{0\dots0}$, so at that input the dependence set is itself a correct slice, the
smallest one that contains every dependence. It makes up \numDepFloor{} of a circuit on average. The QDG slice keeps only
\numLcExcess{} percentage points more and equals it on \numLcExact{} of the \numNPairs{} pairs, while the QPDG keeps
\numYzExcess{} points more, \numYzLcExcessRatio{} times as much code that cannot change the criterion. The gap grows
with width. At \numWElevenLo{} to \numWElevenHi{} qubits the QDG slice keeps \numWElevenLc{} of the circuit and on
corpus~L \numCorpusLSizeLc{}, while the QPDG stays near \numCorpusLSizeYz{}, because it links the operations of a whole
coupling component. Commutation pruning falls below the dependence set, to \numWElevenLcq{} at the same widths,
because it is sound for fixed angles but not complete. It can drop an operation where a fault would change the
criterion, as the identity gate of line~38 in Fig.~\ref{fig:bv} shows, which is why completeness, not size, is what a
developer needs from a slice.

\subsection{RQ2: How much precision beyond the QDG slice is safe?}

\begin{figure}[t]
  \centering
  \includegraphics[width=\columnwidth]{figures/fig_robust.pdf}
  \caption{Share of pairs whose slice, computed at $\ket{0\dots0}$, stays correct in all \numRobDraws{} draws of each
  input regime, for RQ2. Commutation is shown on the \numLcqApplicable{} pairs where Qiskit's pass applies,
  state-based pruning on the \numHThreePrunedN{} pairs it prunes, and angle rebinding on the pairs of circuits with
  rotation angles.}
  \label{fig:robust}
\end{figure}

Both refinements make slices smaller, but only commutation pruning keeps them correct for fixed angles. State-based
pruning shrinks the QDG slice by \numHTwoDiff{} percentage points of the circuit (95\% CI \numHTwoCILow{} to
\numHTwoCIHigh{}) and is smaller on \numHTwoSmaller{} pairs (H2, Wilcoxon, Holm \numHTwoPHolm{}).
Fig.~\ref{fig:robust} shows what that costs. Of the pairs it prunes, the slice stays correct for all random product
inputs on only \numHThreePrunedRate{}, and over all pairs on \numHThreeRate{} (Wilson CI \numHThreeCILow{} to
\numHThreeCIHigh{}, \numCcHThree{} with families as clusters), far below the pre-registered threshold of 95\%, so H3
fails. A state-based slice is a property of one input, not of the program.

Commutation pruning is smaller than the QDG slice on \numLcqSmallerThanLc{} of \numLcqApplicable{} pairs, stays
correct on every random basis and product input, and fails only when the angles change, on \numRobLcqParamFail{} of
\numRobLcqParamN{} pairs of variational circuits. Its symbolic variant, at a mean size of \numLcqsSize{} against
\numLcqSize{}, stays correct on every pair in every regime. Its recall of \numLcqRecall{} shows that it misses
dependences, since a different gate in place of one it drops could change the criterion.

Under depolarizing noise at two-qubit error rates of $10^{-3}$ to $3\cdot10^{-2}$, a pre-registered descriptive
analysis, the criterion's states under the QDG slice and the whole circuit differ by a trace distance of at most
\numNoiseLcCrossExactMax{} where the simulation is exact, rounding error as Theorem~\ref{thm:sound} predicts, while the
state-based slice differs in \numNoiseQCrossExactNonzero{} of \numNoiseQCrossExactN{} exact runs, by up to
\numNoiseQCrossExactMax{}, because it removes noisy gates that the program executes~\cite{preskill2018nisq}.

\subsection{RQ3: What does each rule do to fault localization?}

\begin{table}[t]
  \centering
  \caption{Faults and slices on the \numRqThreeUnits{} seeded-fault units, for RQ3}
  \label{tab:rq3}
  \footnotesize\setlength{\tabcolsep}{3pt}\renewcommand{\arraystretch}{1.12}
  \newcommand{\sep}{\arrayrulecolor{black!25}\hline\arrayrulecolor{black}}
  \begin{tabular}{@{}l r r r r r@{}}
    \toprule
    \textbf{Slice} & \textbf{Contains} & \textbf{Size} & \textbf{Random} & \textbf{Gain} & \textbf{EXAM} \\
    \midrule
    Whole circuit & 100\% & 100\% & -- & -- & \numMbflAllExamAll{} \\
    \midrule
    \rowcolor{blue!7}\textbf{QDG} & \numRqThreeContainLc{} & \numRqThreeSizeLc{} & \numSizeCtlLcRandom{} & \numSizeCtlLcGain{} & \numMbflAllExamLc{} \\ \sep
    QPDG & \numRqThreeContainYz{} & \numRqThreeSizeYz{} & \numSizeCtlYzRandom{} & \numSizeCtlYzGain{} & \numMbflAllExamYz{} \\ \sep
    Classical PDG & \numRqThreeContainCpdg{} & \numRqThreeSizeCpdg{} & \numSizeCtlCpdgRandom{} & \numSizeCtlCpdgGain{} & \numMbflAllExamCpdg{} \\ \sep
    Def-use & \numRqThreeContainDu{} & \numRqThreeSizeDu{} & \numSizeCtlDuRandom{} & \numSizeCtlDuGain{} & \numMbflAllExamDu{} \\ \sep
    State-based & \numRqThreeContainQ{} & \numRqThreeSizeQ{} & \numSizeCtlQRandom{} & \numSizeCtlQGain{} & \numMbflAllExamQ{} \\
    \bottomrule
  \end{tabular}\\[3pt]
  \begin{minipage}{\columnwidth}\footnotesize
  Contains is the share of units whose slice of the faulty circuit contains the fault site, size the mean share of the
  circuit, Random the containment of a random set of the slice's size, and Gain the difference in percentage points.
  EXAM is the mean score of mutation-based fault localization of Section~\ref{sec:design}, lower being better, on
  the \numMbflAllN{} sampled units whose fault every slice contains.
  \end{minipage}
\end{table}

Table~\ref{tab:rq3} shows that a slice that contains the fault more often is also larger, so containment has to be
read against a random set of the slice's size. The QDG slice contains the fault in \numRqThreeContainLc{} of the units
and the QPDG slice in \numRqThreeContainYz{}, but a random set of the QPDG's size already contains it in
\numSizeCtlYzRandom{}, so per unit of size the QDG locates faults better, \numSizeCtlLcGain{} points above random
against \numSizeCtlYzGain{} for the QPDG. Against
a random set of its own size, every rule contains the fault more often, the QDG by \numSizeCtlLcGain{} points (family CI
\numSizeCtlLcGainCi{}), so containment is not only a matter of size. All \numLcMissTotal{} units whose fault the QDG
slice misses are omission faults, a removed control (\numLcMissCrm{}), a re-targeted operand (\numLcMissQsw{}), a
removed preparation of an entangling gate (\numLcMissEnt{}) or a deleted gate (\numLcMissDel{}). Each deletes the
interaction that the slice of the faulty program would follow, so the place of the missing operation is not in the
slice. The QPDG contains \numLcMissYzContained{} of them only because its closure still links operations that the faulty
program no longer connects.

State-based pruning lowers containment by \numHFourDiffAbs{} percentage points against the QDG slice (95\% CI
\numHFourCILow{} to \numHFourCIHigh{}, \numCcHFour{} with families as clusters), so H4, which allowed a loss of at most
five points, fails. In \numCQFour{} units the QDG slice contains the fault and the state-based slice does not. Pruning
removes an operation when the criterion is uncorrelated with its qubits at the input, so a fault that breaks a
correlation can make the faulty operation look irrelevant to the single criterion qubit, and a joint criterion over all
wrong qubits recovers \numCQFourRecovered{} of these units. Where both slices contain the fault, the state-based slice
is smaller by \numHFiveSizeDiff{} operations on average, so H5 holds, but its EXAM score is higher by
\numHFiveExamQHigher{}.

Slices do not improve mutation-based fault localization. On the \numMbflAllN{} units whose fault every slice contains,
ranking within any slice reaches the fault after \numMbflAllExaminedMin{} to \numMbflAllExaminedMax{} operations on
average, as ranking over the whole circuit does, and since EXAM divides by the number of operations ranked, the whole
circuit scores lowest.

On the \numLinesUnits{} real Bugs4Q units whose faulty lines map to the program, the QDG line slices keep every faulty
line and \numLinesLcMean{} of \numLinesFullMean{} lines on average, so on these small programs slicing barely reduces
what there is to read. Def-use, which equals the classical PDG on these unitary programs, loses the faulty line in
\numLinesDuMissed{} of the units.

\paragraph*{Language models}
Asked to list the operations that can affect a qubit, Claude Sonnet~5.5 and Haiku~5.5 recall \numTOneSonnetRecall{}
and \numTOneHaikuRecall{} of the dependences and both Qwen models \numTOneQwenMidRecall{}, where the QDG slice recalls
them all, and Sonnet recalls only \numTOneSonnetKick{} of those through kickback. Asked to name the faulty operation,
no model does better with any slice than with the whole circuit by more than a fraction of a point, and Sonnet has
the fault in its top three for \numTTwoSonnetQ{} of the units with the state-based slice, \numTTwoSonnetQVsLcLoss{}
points below the QDG slice (95\% CI \numTTwoSonnetQVsLcCI{}, exact McNemar \numTTwoSonnetQVsLcP{}). A model thus misses
dependences that a slicer finds by rule, yet a slice does not help it localize a fault, and a pruned slice hurts it.

\subsection{RQ4: What does it cost?}

The QDG slice costs at most \numCostLcMaxCritMs{}~ms per criterion at every depth.
State-based pruning, which simulates the circuit for every candidate removal, exceeds the limit of \numCostWallLimit{}~s
for \numCostDepthEightyOverLimit{} of the \numCostDepthCircuits{} circuits at the second-largest depth and for all
\numCostDepthOneSixtyOverLimit{} at the largest, grows exponentially with width, and completes on
\numCorpusLCompleted{} of the \numCorpusLCircuits{} circuits of corpus~L. Building the QPDG takes
\numCostBuildDepthForty{}~s at depth \numCostMeasuredDepthForty{} and \numCostBuildDepthOneSixty{}~s at depth
\numCostMeasuredDepthOneSixty{} in the median, against milliseconds for the QDG slice, because its closure links every
pair of operations on a coupling component while the QDG needs only the operations on each qubit.


\section{Implications for Comprehension Tools}\label{sec:implications}

\paragraph*{What a slicer must do}
A slicer must let every gate write all its qubits, so a representation that reads the control as unchanged, as the
QCPG and order tables built for pattern detection do, should not be read as a dependence graph.
A slicer must also follow measured bits, without which def-use is wrong on \numDuDynamicUnsound{} of the
\numNDynamicPairs{} dynamic pairs, and should not add entanglement edges, which add \numYzLcExcessRatio{} times as much
code that cannot matter and take seconds to build where the QDG takes milliseconds. Since one loop iteration can put a
line in a slice, it should show the iterations that do.

\paragraph*{What a developer should not trust}
A slice cannot be validated by its outcomes, as \numFdCpdgZMatchPct{} of the wrong classical slices and the Bugs4Q
example show, so a tool that checks slices should compare reduced states, or at least joint outcomes in more than one
basis. Pruned slices
answer a narrower question. A commutation-pruned slice is right for one observable at the circuit's angles. A state-based
slice is right for one input, stays correct for random inputs on only \numHThreeRate{} of the pairs, and contains the
fault \numHFourDiffAbs{} points less often. A tool should label either as such and show the QDG slice by default.

\paragraph*{What a slice cannot show}
Every unit whose fault the QDG slice missed is an omission fault, a removed control, gate or preparation, or a gate
moved to another qubit. No slice contains an operation the program lacks, so when a qubit's slice leaves out a qubit
the developer expected, that absence is the clue, and listing the qubits each slice contains makes it visible.

\paragraph*{What remains open}
Slices narrowed where to look but improved neither fault ranking nor a language model's localization, while the models
missed dependences that the rule finds. Whether slices shorten the manual reading that developers report does not
scale~\cite{zappin2026challenges} needs a study with developers, which the semantics and measured rules of this paper
make possible to design.


\section{Threats to Validity}\label{sec:threats}

\paragraph*{Construct validity}
The dependence set is measured with \numInfluencePerturbations{} random perturbations per operation at
\numInfluenceInputs{} inputs, so it can miss true dependences but never adds false ones, and the QDG's distance to it
may be smaller than reported. The dependence set and the wrong slices of every rule were checked at $\ket{0\dots0}$,
which RQ2 extends to random inputs. A slice preserves its criterion at a squared fidelity of at least $1 - 10^{-6}$,
far above the oracle's error. Each criterion is a single qubit, and Proposition~\ref{prop:joint} extends the QDG to
sets but not the pruning rules. Fault containment measures where a developer would look, not whether a developer finds
the fault. The QPDG is reimplemented from prose, with every reading logged before any run and three variants
bracketing the plausible ones. Qiskit's pass applies to \numLcqApplicable{} of the \numNPairs{} pairs, those without
mid-circuit measurements.

\paragraph*{Internal validity}
The oracle agrees with Qiskit's simulation to within \numOracleErr{}, a catalog of \numCatEntries{} hand-derived cases
checks the ground truth and every slicer before the runs, and the nesting of the rules, which holds on every pair,
would expose a slicer that drops or adds operations. The soundness proof is mechanized for an abstract model of local
operations, and that quantum operations satisfy the model is argued on paper. The outcomes of H1 and H2 were known
when the protocol was written, and both hold by construction. Every later change is a dated deviation, and every
result outside the five hypotheses is exploratory. The Claude models accept only their default sampling, so their
answers can differ between calls.

\paragraph*{External validity}
The circuits with an exact ground truth have at most \numSizeAMaxQubits{} qubits in corpus~A and 20 in corpus~L, and
\numNDynamicPairs{} pairs come from dynamic circuits. The proofs hold for every circuit, but the measured distances
between the rules may change with width and feedback. Benchmark circuits are not programs written by developers, the
\numLinesUnits{} Bugs4Q units are the only real faults, and seeded faults may not match real ones in their mix. The
gate check uses one angle per gate, the noise analysis depolarizing noise only, and loops that run on the device are
not studied.


\section{Conclusion}\label{sec:conclusion}

A quantum slice helps a developer only if it never hides code that can change the qubits asked about. We defined that
requirement as preserving the qubits' joint state for every input while keeping every semantic dependence, and showed
that the write-all rule meets it. Its slice is the backward light cone extended along classical bits, needs only
classical data, classical control and quantum data dependence, and is computed in one pass. A pre-registered study
measured what the other rules lose. Reading a control as unchanged gives wrong slices that often reproduce the
program's outcomes, entanglement edges add only code that cannot matter, and pruning trades dependences for size and
holds for one observable or input. The QDG slice comes within a few points of the smallest complete slice and costs
milliseconds. Two questions remain, whether these slices shorten the manual reading that developers report does not
scale, which needs a study with developers, and how to compute them statically on the source developers write, at
the scale of real programs and with loops that run on the device.


\bibliographystyle{IEEEtran}
\bibliography{../references}
\end{document}
